There are many sources (such as \cite{gps1995,gurski2013tricatcoh,yau2cat2021}) that discuss the strictification trifunctor $\st \colon \BiCat \rightarrow \nCat{2}_{\st,\ps}$ and its supposed right \emph{triadjoint} $i\colon \nCat{2}_{st,ps} \rightarrow \BiCat$. Here, $\nCat{2}_{\st,\ps}$ is also known as $\Gray$, and is the tricategory of strict $2$-categories, strict $2$-functors, pseudonatural transformations, and modifications. We aim to expand upon these descriptions by providing a more thorough account of this functor using concrete diagrams to fully specify the relevant trifunctorial data. Then, we shall briefly describe its usage to encode axioms of bicategorical structures using pasting diagrams. 

\subsection{Mapping of Bicategories}

Each bicategory $B \in \BiCat$ maps to a strict $2$-category called $\st B$, which is essentially determined by a pair of maps $e \colon i\circ \st B \rightarrow B$ and $s \colon B \rightarrow i \circ \st B$. For simplicity, we shall drop the $i$, unless required for better clarity. We first define the action of $e$ on the objects and $1$-cells of $\st B$. 

\begin{definition}
	The map $e$ is defined on $0$-cells and $1$-cells as follows:
	\begin{description}
		\item[$0$-cells -] for each $b \in \st B$, $e(b) = b \in B$,
		\item[$1$-cells -] for each non-empty string $(f_n,f_{n-1},\cdots,f_1)$ representing a $1$-cell in $\st B$, written right to left, $e$ is defined recursively as shown below:
		% https://q.uiver.app/#q=WzAsMjEsWzAsMCwiYiJdLFsxLDAsImJfMSJdLFsyLDAsImJfMiJdLFszLDAsIlxcY2RvdHMiXSxbNCwwLCJiX24iXSxbNSwwLCJcXG1hcHN0b157ZX0iXSxbNiwwLCJiIl0sWzgsMCwiYl9uIl0sWzEsMSwiYl8xIl0sWzIsMSwiYl8yIl0sWzMsMSwiXFxjZG90cyJdLFs0LDEsImJfbiJdLFs1LDEsIlxcbWFwc3RvXntlfSJdLFs2LDEsImJfMSJdLFs4LDEsImJfbiJdLFs0LDMsImJfbiJdLFs0LDIsIlxcdmRvdHMiXSxbMywzLCJiX3tuLTF9Il0sWzUsMywiXFxtYXBzdG9ee2V9Il0sWzYsMywiYl97bi0xfSJdLFs4LDMsImJfbiJdLFswLDEsImZfMSJdLFsxLDIsImZfMiJdLFsyLDMsImZfMyJdLFszLDRdLFs2LDcsImUoZl9uLFxcY2RvdHMsZl8yKWZfMSJdLFs4LDksImZfMiJdLFs5LDEwLCJmXzMiXSxbMTAsMTFdLFsxMywxNCwiZShmX24sXFxjZG90cyxmXzMpZl8yIl0sWzE3LDE1LCJmX24iXSxbMTksMjAsImZfbiJdXQ==
		\[\begin{tikzcd}[ampersand replacement=\&,cramped,row sep=small]
			b \& {b_1} \& {b_2} \& \cdots \& {b_n} \& {\mapsto^{e}} \& b \&\& {b_n} \\
			\& {b_1} \& {b_2} \& \cdots \& {b_n} \& {\mapsto^{e}} \& {b_1} \&\& {b_n} \\
			\&\&\&\& \vdots \\
			\&\&\& {b_{n-1}} \& {b_n} \& {\mapsto^{e}} \& {b_{n-1}} \&\& {b_n}
			\arrow["{f_1}", from=1-1, to=1-2]
			\arrow["{f_2}", from=1-2, to=1-3]
			\arrow["{f_3}", from=1-3, to=1-4]
			\arrow[from=1-4, to=1-5]
			\arrow["{e(f_n,\cdots,f_2)f_1}", from=1-7, to=1-9]
			\arrow["{f_2}", from=2-2, to=2-3]
			\arrow["{f_3}", from=2-3, to=2-4]
			\arrow[from=2-4, to=2-5]
			\arrow["{e(f_n,\cdots,f_3)f_2}", from=2-7, to=2-9]
			\arrow["{f_n}", from=4-4, to=4-5]
			\arrow["{f_n}", from=4-7, to=4-9]
		\end{tikzcd}\cdpunct[2pt]{,}\]
		and for the empty string, it is defined seperately by setting $(b \xrightarrow{\varepsilon_b} b) \mapsto^{e} (1_b \colon b \rightarrow b)$.
	\end{description}
\end{definition}

Having partially defined $e$, we can now fully specify the $2$-category $\st B$ in terms of its $2$-cells, vertical composition, and horiztonal composition.

\begin{definition}
	The $2$-category $\st B$ has the following data:
	\begin{description}
		\item[Objects -] exactly the objects of $B$,
		\item[$1$-cells -] free strings of arrows $b \xrightarrow{f_1} b_1 \xrightarrow{f_2} b_2 \xrightarrow{f_2} b_3 \xrightarrow{f_3} \cdots $, where the empty string is allowed and is labeled as $b \xrightarrow{\varepsilon_{b}} b$,
		\item[Identity $1$-cells -] for each $b$, the empty string $b \xrightarrow{\varepsilon_{b}} b$ is the horizontal identity
		\item[Hcomp $1$-cells -] two strings of arrows are composed by concatenation of strings,
		\item[$2$-cells -] the $2$-cells between two strings $\vec{f}$ and $\vec{g}$ are in bijection with the $2$-cells $\alpha \colon e(\vec{f}) \Rightarrow e(\vec{g})$, % https://q.uiver.app/#q=WzAsMTYsWzIsMywiYl8yIl0sWzEsMCwiQiJdLFs1LDAsIlxcc3QgQiJdLFs0LDIsImIiXSxbNSwyLCJiXzEiXSxbNiwyLCJiXzIiXSxbNywyLCJcXGNkb3RzIl0sWzgsMiwiYl9uIl0sWzQsNSwiYiJdLFs1LDUsImInXzEiXSxbNiw1LCJiJ18yIl0sWzcsNSwiXFxjZG90cyJdLFs4LDUsImInX3ttfT1iX24iXSxbMCwzLCJiXzEiXSxbMiw0LCJiXzIiXSxbMCw0LCJiXzEiXSxbMyw0LCJmXzEiXSxbNCw1LCJmXzIiXSxbNSw2LCJmXzMiXSxbNiw3XSxbOCw5LCJnXzEiLDJdLFs5LDEwLCJnXzIiLDJdLFsxMCwxMSwiZ18zIiwyXSxbMTEsMTJdLFsxMywwLCJlKFxcdmVje2Z9KSIsMCx7ImN1cnZlIjotMn1dLFsxNSwxNCwiZShcXHZlY3tnfSkiLDIseyJjdXJ2ZSI6Mn1dLFsyLDEsImUiLDJdLFs1LDEwLCIoXFxhbHBoYSA6IFxcdmVje2Z9IFxcUmlnaHRhcnJvdyBcXHZlY3tnfSBcXGluIEIpIiwwLHsic2hvcnRlbiI6eyJzb3VyY2UiOjIwLCJ0YXJnZXQiOjIwfSwibGV2ZWwiOjJ9XSxbMjQsMjUsIlxcYWxwaGEiLDAseyJzaG9ydGVuIjp7InNvdXJjZSI6MjAsInRhcmdldCI6MjB9fV1d
		\[\begin{tikzcd}[ampersand replacement=\&]
			\& B \&\&\&\& {\st B} \&\&\& \\
			\\
			\&\&\&\& b \& {b_1} \& {b_2} \& \cdots \& {b_n} \\
			{b_1} \&\& {b_2} \\
			{b_1} \&\& {b_2} \\
			\&\&\&\& b \& {b'_1} \& {b'_2} \& \cdots \& {b'_{m}=b_n}
			\arrow["e"', from=1-6, to=1-2]
			\arrow["{f_1}", from=3-5, to=3-6]
			\arrow["{f_2}", from=3-6, to=3-7]
			\arrow["{f_3}", from=3-7, to=3-8]
			\arrow["{(\alpha \colon \vec{f} \Rightarrow \vec{g} \in B)}", between={0.2}{0.8}, Rightarrow, nfold, from=3-7, to=6-7]
			\arrow[from=3-8, to=3-9]
			\arrow[""{name=0, anchor=center, inner sep=0}, "{e(\vec{f})}", curve={height=-12pt}, from=4-1, to=4-3]
			\arrow[""{name=1, anchor=center, inner sep=0}, "{e(\vec{g})}"', curve={height=12pt}, from=5-1, to=5-3]
			\arrow["{g_1}"', from=6-5, to=6-6]
			\arrow["{g_2}"', from=6-6, to=6-7]
			\arrow["{g_3}"', from=6-7, to=6-8]
			\arrow[from=6-8, to=6-9]
			\arrow["\alpha", between={0.2}{0.8}, Rightarrow, nfold, from=0, to=1]
		\end{tikzcd}\cdpunct[8pt]{,}\]
		this bijection of $2$-cells also defines the action of $e$ on the $2$-cells of $\st B$,
		\item[Identity $2$-cells -] for each string $\vec{f}$, the unique $2$-cell $1_{e({\vec{f}})} \colon e(\vec{f}) \Rightarrow e(\vec{f})$ is the identity,
		\item[Vcomp $2$-cells -] the vertical composition of $2$-cells is directly lifted from the map $e$ as shown below
		% https://q.uiver.app/#q=WzAsMjMsWzQsMywiYl8yIl0sWzMsMCwiQiJdLFs3LDAsIlxcc3QgQiJdLFs2LDIsImIiXSxbNywyLCJiXzEiXSxbOCwyLCJiXzIiXSxbOSwyLCJcXGNkb3RzIl0sWzEwLDIsImJfbiJdLFs1LDUsImIiXSxbNiw1LCJiJ18xIl0sWzcsNSwiYidfMiJdLFs4LDUsIlxcY2RvdHMiXSxbOSw1LCJiJ197bX09Yl9uIl0sWzIsMywiYl8xIl0sWzIsNSwiYl8yIl0sWzAsNSwiYl8xIl0sWzIsNywiYl8xIl0sWzQsNywiYl8yIl0sWzYsOCwiYiJdLFsxMCw4LCJiJydfe3B9ID0gYl9uIl0sWzcsOCwiYicnXzEiXSxbOCw4LCJiJydfMiJdLFs5LDgsIlxcY2RvdHMiXSxbMyw0LCJmXzEiXSxbNCw1LCJmXzIiXSxbNSw2LCJmXzMiXSxbNiw3XSxbOCw5LCJnXzEiLDJdLFs5LDEwLCJnXzIiLDJdLFsxMCwxMSwiZ18zIiwyXSxbMTEsMTJdLFsxMywwLCJlKFxcdmVje2Z9KSIsMCx7ImN1cnZlIjotMn1dLFsxNSwxNCwiZShcXHZlY3tnfSkiLDJdLFsyLDEsImUiLDJdLFs1LDEwLCIoXFxhbHBoYSA6IFxcdmVje2Z9IFxcUmlnaHRhcnJvdyBcXHZlY3tnfSBcXGluIEIpIiwwLHsic2hvcnRlbiI6eyJzb3VyY2UiOjIwLCJ0YXJnZXQiOjIwfSwibGV2ZWwiOjJ9XSxbMTYsMTcsImUoXFx2ZWN7aH0pIiwyLHsiY3VydmUiOjJ9XSxbMjAsMjEsImhfMiIsMl0sWzIxLDIyLCJoXzMiLDJdLFsyMiwxOV0sWzE4LDIwLCJoXzEiLDJdLFsxMCwyMSwiKFxcYmV0YSA6IFxcdmVje2d9IFxcUmlnaHRhcnJvdyBcXHZlY3tofSBcXGluIEIpIiwwLHsic2hvcnRlbiI6eyJzb3VyY2UiOjIwLCJ0YXJnZXQiOjIwfSwibGV2ZWwiOjJ9XSxbMzEsMzIsIlxcYWxwaGEiLDAseyJzaG9ydGVuIjp7InNvdXJjZSI6MjAsInRhcmdldCI6MjB9fV0sWzMyLDM1LCJcXGJldGEiLDAseyJzaG9ydGVuIjp7InNvdXJjZSI6MjAsInRhcmdldCI6MjB9fV0sWzMxLDM1LCJcXGJldGFcXGFscGhhIiwwLHsic2hvcnRlbiI6eyJzb3VyY2UiOjIwLCJ0YXJnZXQiOjIwfX1dLFsyNiwzOCwiXFxiZXRhXFxhbHBoYSIsMCx7Im9mZnNldCI6LTUsImN1cnZlIjotNSwic2hvcnRlbiI6eyJzb3VyY2UiOjIwLCJ0YXJnZXQiOjIwfX1dXQ==
		\[\begin{tikzcd}[ampersand replacement=\&,column sep=small]
			\&\&\& B \&\&\&\& {\st B} \&\&\& \\
			\\
			\&\&\&\&\&\& b \& {b_1} \& {b_2} \& \cdots \& {b_n} \\
			\&\& {b_1} \&\& {b_2} \\
			\\
			{b_1} \&\& {b_2} \&\&\& b \& {b'_1} \& {b'_2} \& \cdots \& {b'_{m}=b_n} \\
			\\
			\&\& {b_1} \&\& {b_2} \\
			\&\&\&\&\&\& b \& {b''_1} \& {b''_2} \& \cdots \& {b''_{p} = b_n}
			\arrow["e"', from=1-8, to=1-4]
			\arrow["{f_1}", from=3-7, to=3-8]
			\arrow["{f_2}", from=3-8, to=3-9]
			\arrow["{f_3}", from=3-9, to=3-10]
			\arrow["{(\alpha \colon \vec{f} \Rightarrow \vec{g} \in B)}", between={0.2}{0.8}, Rightarrow, nfold, from=3-9, to=6-8]
			\arrow[""{name=0, anchor=center, inner sep=0}, from=3-10, to=3-11]
			\arrow[""{name=1, anchor=center, inner sep=0}, "{e(\vec{f})}", curve={height=-12pt}, from=4-3, to=4-5]
			\arrow[""{name=2, anchor=center, inner sep=0}, "{e(\vec{g})}"', from=6-1, to=6-3]
			\arrow["{g_1}"', from=6-6, to=6-7]
			\arrow["{g_2}"', from=6-7, to=6-8]
			\arrow["{g_3}"', from=6-8, to=6-9]
			\arrow["{(\beta \colon \vec{g} \Rightarrow \vec{h} \in B)}", between={0.2}{0.8}, Rightarrow, nfold, from=6-8, to=9-9]
			\arrow[from=6-9, to=6-10]
			\arrow[""{name=3, anchor=center, inner sep=0}, "{e(\vec{h})}"', curve={height=12pt}, from=8-3, to=8-5]
			\arrow["{h_1}"', from=9-7, to=9-8]
			\arrow["{h_2}"', from=9-8, to=9-9]
			\arrow["{h_3}"', from=9-9, to=9-10]
			\arrow[""{name=4, anchor=center, inner sep=0}, from=9-10, to=9-11]
			\arrow["{\beta\alpha}", shift left=5, curve={height=-30pt}, between={0.2}{0.8}, Rightarrow, nfold, from=0, to=4]
			\arrow["\alpha", between={0.2}{0.8}, Rightarrow, nfold, from=1, to=2]
			\arrow["{\beta\alpha}", between={0.2}{0.8}, Rightarrow, nfold, from=1, to=3]
			\arrow["\beta", between={0.2}{0.8}, Rightarrow, nfold, from=2, to=3]
		\end{tikzcd}\]
		\item[Hcomp $2$-cells -] the horiztonal composite of two $2$-cells $\alpha \colon e(\vec{f}) \Rightarrow e(\vec{g})$ and $\beta \colon e(\vec{h}) \Rightarrow e(\vec{k})$ is given as follows:
		% https://q.uiver.app/#q=WzAsMzUsWzAsMCwiYiJdLFsxLDAsImJfMSJdLFsyLDAsIlxcY2RvdHMiXSxbMywwLCJiX249YyJdLFswLDIsImIiXSxbMiwyLCJcXGNkb3RzIl0sWzEsMiwiYidfMSJdLFszLDIsImInX3ttfT1iX249YyJdLFs0LDAsImNfMSJdLFs0LDIsImMnXzEiXSxbNSwwLCJcXGNkb3RzIl0sWzYsMCwiY197cH09ZCJdLFs1LDIsIlxcY2RvdHMiXSxbNiwyLCJjJ19xID0gY197cH09ZCJdLFszLDMsIlxcdmVxdWl2Il0sWzAsNCwiYiJdLFsxLDQsImJfMSJdLFsyLDQsIlxcY2RvdHMiXSxbMyw0LCJiX249YyJdLFs0LDQsImNfMSJdLFs1LDQsIlxcY2RvdHMiXSxbNiw0LCJjX3A9ZCJdLFswLDYsImIiXSxbMyw2LCJjIl0sWzYsNiwiZCJdLFswLDgsImIiXSxbMyw4LCJjIl0sWzYsOCwiZCJdLFswLDEwLCJiIl0sWzEsMTAsImInXzEiXSxbMiwxMCwiXFxjZG90cyJdLFszLDEwLCJiJ19tPWJfbj1jIl0sWzQsMTAsImMnXzEiXSxbNSwxMCwiXFxjZG90cyJdLFs2LDEwLCJjJ19xPWNfcD1kIl0sWzEsMiwiZl8yIl0sWzIsM10sWzYsNSwiZ18yIiwyXSxbNSw3XSxbMCwxLCJmXzEiXSxbNCw2LCJnXzEiLDJdLFszLDgsImhfMSJdLFs3LDksImtfMSIsMl0sWzgsMTAsImhfMiJdLFsxMCwxMV0sWzksMTIsImtfMiIsMl0sWzEyLDEzXSxbMjgsMjksImdfMSIsMl0sWzI5LDMwLCJnXzIiLDJdLFszMCwzMV0sWzMxLDMyLCJrXzEiLDJdLFszMiwzMywia18yIiwyXSxbMzMsMzRdLFsyNSwyNiwiZShcXHZlY3tnfSkiLDFdLFsyNiwyNywiZShcXHZlY3trfSkiLDFdLFsxNSwxNiwiZl8xIl0sWzE2LDE3LCJmXzIiXSxbMTcsMThdLFsxOCwxOSwiaF8xIl0sWzE5LDIwLCJoXzIiXSxbMjAsMjFdLFsyMiwyMywiZShcXHZlY3tmfSkiLDFdLFsyMywyNCwiZShcXHZlY3tofSkiLDFdLFsxOCwyMywiZShcXHZlY3tofSxcXHZlY3tmfSkgXFxSaWdodGFycm93IGUoXFx2ZWN7aH0pIGUoXFx2ZWN7Zn0pIiwwLHsic2hvcnRlbiI6eyJzb3VyY2UiOjIwLCJ0YXJnZXQiOjIwfSwibGV2ZWwiOjJ9XSxbMjYsMzEsIiBlKFxcdmVje2t9KWUoXFx2ZWN7Z30pXFxSaWdodGFycm93IGUoXFx2ZWN7a30sXFx2ZWN7Z30pIiwwLHsic2hvcnRlbiI6eyJzb3VyY2UiOjIwLCJ0YXJnZXQiOjIwfSwibGV2ZWwiOjJ9XSxbMjMsMjYsIlxcYmV0YSBcXHN0YXIgXFxhbHBoYSIsMSx7InNob3J0ZW4iOnsic291cmNlIjoyMCwidGFyZ2V0IjoyMH0sImxldmVsIjoyfV0sWzM1LDM3LCIoXFxhbHBoYSA6IFxcdmVje2Z9IFxcUmlnaHRhcnJvdyBcXHZlY3tnfSBcXGluIEIpIiwyLHsic2hvcnRlbiI6eyJzb3VyY2UiOjIwLCJ0YXJnZXQiOjIwfX1dLFs0Myw0NSwiKFxcYmV0YSA6IFxcdmVje2h9IFxcUmlnaHRhcnJvdyBcXHZlY3trfSBcXGluIEIpIiwwLHsic2hvcnRlbiI6eyJzb3VyY2UiOjIwLCJ0YXJnZXQiOjIwfX1dLFs2Niw2NywiXFxzdGFyIiwxLHsic2hvcnRlbiI6eyJzb3VyY2UiOjIwLCJ0YXJnZXQiOjIwfSwic3R5bGUiOnsiYm9keSI6eyJuYW1lIjoibm9uZSJ9LCJoZWFkIjp7Im5hbWUiOiJub25lIn19fV1d
		\[\begin{tikzcd}[ampersand replacement=\&]
			b \& {b_1} \& \cdots \& {b_n=c} \& {c_1} \& \cdots \& {c_{p}=d} \\
			\\
			b \& {b'_1} \& \cdots \& {b'_{m}=b_n=c} \& {c'_1} \& \cdots \& {c'_q = c_{p}=d} \\
			\&\&\& \vequiv \\
			b \& {b_1} \& \cdots \& {b_n=c} \& {c_1} \& \cdots \& {c_p=d} \\
			\\
			b \&\&\& c \&\&\& d \\
			\\
			b \&\&\& c \&\&\& d \\
			\\
			b \& {b'_1} \& \cdots \& {b'_m=b_n=c} \& {c'_1} \& \cdots \& {c'_q=c_p=d}
			\arrow["{f_1}", from=1-1, to=1-2]
			\arrow[""{name=0, anchor=center, inner sep=0}, "{f_2}", from=1-2, to=1-3]
			\arrow[from=1-3, to=1-4]
			\arrow["{h_1}", from=1-4, to=1-5]
			\arrow[""{name=1, anchor=center, inner sep=0}, "{h_2}", from=1-5, to=1-6]
			\arrow[from=1-6, to=1-7]
			\arrow["{g_1}"', from=3-1, to=3-2]
			\arrow[""{name=2, anchor=center, inner sep=0}, "{g_2}"', from=3-2, to=3-3]
			\arrow[from=3-3, to=3-4]
			\arrow["{k_1}"', from=3-4, to=3-5]
			\arrow[""{name=3, anchor=center, inner sep=0}, "{k_2}"', from=3-5, to=3-6]
			\arrow[from=3-6, to=3-7]
			\arrow["{f_1}", from=5-1, to=5-2]
			\arrow["{f_2}", from=5-2, to=5-3]
			\arrow[from=5-3, to=5-4]
			\arrow["{h_1}", from=5-4, to=5-5]
			\arrow["{e(\vec{h},\vec{f}) \Rightarrow e(\vec{h}) e(\vec{f})}", between={0.2}{0.8}, Rightarrow, nfold, from=5-4, to=7-4]
			\arrow["{h_2}", from=5-5, to=5-6]
			\arrow[from=5-6, to=5-7]
			\arrow["{e(\vec{f})}"{description}, from=7-1, to=7-4]
			\arrow["{e(\vec{h})}"{description}, from=7-4, to=7-7]
			\arrow["{\beta \star \alpha}"{description}, between={0.2}{0.8}, Rightarrow, nfold, from=7-4, to=9-4]
			\arrow["{e(\vec{g})}"{description}, from=9-1, to=9-4]
			\arrow["{e(\vec{k})}"{description}, from=9-4, to=9-7]
			\arrow["{ e(\vec{k})e(\vec{g})\Rightarrow e(\vec{k},\vec{g})}", between={0.2}{0.8}, Rightarrow, nfold, from=9-4, to=11-4]
			\arrow["{g_1}"', from=11-1, to=11-2]
			\arrow["{g_2}"', from=11-2, to=11-3]
			\arrow[from=11-3, to=11-4]
			\arrow["{k_1}"', from=11-4, to=11-5]
			\arrow["{k_2}"', from=11-5, to=11-6]
			\arrow[from=11-6, to=11-7]
			\arrow[""{name=4, anchor=center, inner sep=0}, "{(\alpha \colon \vec{f} \Rightarrow \vec{g} \in B)}"', between={0.2}{0.8}, Rightarrow, nfold, from=0, to=2]
			\arrow[""{name=5, anchor=center, inner sep=0}, "{(\beta \colon \vec{h} \Rightarrow \vec{k} \in B)}", between={0.2}{0.8}, Rightarrow, nfold, from=1, to=3]
			\arrow["\star"{description}, draw=none, from=4, to=5]
		\end{tikzcd}\cdpunct[8pt]{,}\]
		where the $2$-cells $e(\vec{k})e(\vec{g})\Rightarrow e(\vec{k},\vec{g})$ and $e(\vec{h},\vec{f}) \Rightarrow e(\vec{h}) e(\vec{f})$ are induced uniquely from the coherence constraints, and $\beta \star \alpha$ is the horizontal composite in the bicategory $B$ itself.
	\end{description}
\end{definition}

Now, we wish to specify $e \colon i\circ \st B \rightarrow B$ and $s \colon B \rightarrow i \circ \st B$ as pseudofunctors.

\begin{definition}
	The pseudofunctor $e \colon \st B \rightarrow B$ is defined by the following data:
	\begin{description}
		\item[Objects -] for each $b\in \st B$, the pseudofunctor $e$ acts as $b \mapsto b$,
		\item[$1$-cells -] for each string $\vec{f} \in \st B$, the pseudofunctor $e$ acts as $\vec{f} \mapsto e(\vec{f})$,
		\item[$2$-cells -] each $2$-cell in $\st B$ is already defined as a $2$-cell in $B$,
		\item[Composition constraint - ] the $2$-cell $e(\vec{g}) e(\vec{f}) \Rightarrow e(\vec{g},\vec{f})$ is the unique coherence constraint $2$-cell,
		% https://q.uiver.app/#q=WzAsOSxbMCwyLCJiIl0sWzEsMSwiYyJdLFsyLDIsImQiXSxbMSwwLCJcXHN0IEIiXSxbNSwwLCJCIl0sWzQsMiwiYiJdLFs1LDEsImMiXSxbNiwyLCJkIl0sWzMsMiwiXFxyaWdodHNxdWlnYXJyb3ciXSxbMSwyLCJcXHZlY3tnfSJdLFswLDEsIlxcdmVje2Z9Il0sWzMsNCwiZSJdLFswLDIsIlxcdmVje2d9LFxcdmVje2Z9IiwyXSxbNSw2LCJlKFxcdmVje2Z9KSJdLFs2LDcsImUoXFx2ZWN7Z30pIl0sWzUsNywiZShcXHZlY3tnfSxcXHZlY3tmfSkiLDJdLFsxLDEyLCIiLDAseyJzaG9ydGVuIjp7InNvdXJjZSI6MjAsInRhcmdldCI6MjB9LCJzdHlsZSI6eyJoZWFkIjp7Im5hbWUiOiJub25lIn19fV0sWzYsMTUsIiIsMCx7InNob3J0ZW4iOnsic291cmNlIjoyMCwidGFyZ2V0IjoyMH19XV0=
		\[\begin{tikzcd}[ampersand replacement=\&]
			\& {\st B} \&\&\&\& B \& \\
			\& c \&\&\&\& c \\
			b \&\& d \& \rightsquigarrow \& b \&\& d
			\arrow["e", from=1-2, to=1-6]
			\arrow["{\vec{g}}", from=2-2, to=3-3]
			\arrow["{e(\vec{g})}", from=2-6, to=3-7]
			\arrow["{\vec{f}}", from=3-1, to=2-2]
			\arrow[""{name=0, anchor=center, inner sep=0}, "{\vec{g},\vec{f}}"', from=3-1, to=3-3]
			\arrow["{e(\vec{f})}", from=3-5, to=2-6]
			\arrow[""{name=1, anchor=center, inner sep=0}, "{e(\vec{g},\vec{f})}"', from=3-5, to=3-7]
			\arrow[between={0.2}{0.8}, equals, nfold, from=2-2, to=0]
			\arrow[between={0.2}{0.8}, Rightarrow, nfold, from=2-6, to=1]
		\end{tikzcd}\cdpunct[8pt]{,}\]
		\item[Identity constraint - ] the $2$-cell $1_b \Rightarrow e(b \xrightarrow{\varepsilon_{b}} b)$ induced by the identity $2$-cell at $1_b \colon b \rightarrow b$,
		% https://q.uiver.app/#q=WzAsNyxbMCwyLCJiIl0sWzIsMiwiYiJdLFsxLDAsIlxcc3QgQiJdLFs1LDAsIkIiXSxbNCwyLCJiIl0sWzYsMiwiZCJdLFszLDIsIlxccmlnaHRzcXVpZ2Fycm93Il0sWzIsMywiZSJdLFswLDEsIlxcdmFyZXBzaWxvbl97Yn0iLDJdLFs0LDUsIjFfe2J9PWUoXFx2YXJlcHNpbG9uX3tifSkiLDIseyJjdXJ2ZSI6Mn1dLFs0LDUsIjFfe2J9IiwwLHsiY3VydmUiOi0yfV0sWzEwLDksIjFfezFfYn0iLDAseyJzaG9ydGVuIjp7InNvdXJjZSI6MjAsInRhcmdldCI6MjB9fV1d
		\[\begin{tikzcd}[ampersand replacement=\&]
			\& {\st B} \&\&\&\& B \& \\
			\\
			b \&\& b \& \rightsquigarrow \& b \&\& d
			\arrow["e", from=1-2, to=1-6]
			\arrow["{\varepsilon_{b}}"', from=3-1, to=3-3]
			\arrow[""{name=0, anchor=center, inner sep=0}, "{1_{b}=e(\varepsilon_{b})}"', curve={height=12pt}, from=3-5, to=3-7]
			\arrow[""{name=1, anchor=center, inner sep=0}, "{1_{b}}", curve={height=-12pt}, from=3-5, to=3-7]
			\arrow["{1_{1_b}}", between={0.2}{0.8}, Rightarrow, nfold, from=1, to=0]
		\end{tikzcd}\cdpunct[8pt]{.}\]
	\end{description}
\end{definition}

The pseudofunctor $s \colon B \rightarrow \st B$ is supposed to be a strict section of the pseudofunctor $e \colon \st B \rightarrow B$, that is, $e \circ s = 1_{B}$. We shall specify $s$ with this view in mind.

\begin{definition}
	The pseudofunctor $s \colon B \rightarrow \st B$ is defined as follows:
	\begin{description}
		\item[Objects -] for each $b\in B$, the pseudofunctor $s$ acts as $b \mapsto b$,
		\item[$1$-cells -] for each $f \colon b \rightarrow c \in B$, the pseudofunctor $f$ acts as $f \mapsto (f)$, where $(f)$ is string of length one, for simplicity this will now be labeled as $f$ itself,
		\item[$2$-cells -] the $2$-cell $\alpha \colon f \Rightarrow g$ is mapped precisely to the $2$-cell $\alpha \colon f \Rightarrow g$ under inverse image in $e$,
		\item[Composition constraint -] the $2$-cell $(g,f) \rightarrow gf$ is induced by the identity cell $1_{gf} \colon e(g,f) \Rightarrow e(gf)$,
		% https://q.uiver.app/#q=WzAsOSxbNCwyLCJiIl0sWzUsMSwiYyJdLFs2LDIsImQiXSxbNSwwLCJcXHN0IEIiXSxbMSwwLCJCIl0sWzAsMiwiYiJdLFsxLDEsImMiXSxbMiwyLCJkIl0sWzMsMiwiXFxyaWdodHNxdWlnYXJyb3ciXSxbMSwyLCJnIl0sWzAsMSwiZiJdLFs0LDMsImYiLDJdLFswLDIsImdmIiwyXSxbNSw2LCJmIl0sWzYsNywiZyJdLFs1LDcsImdmIiwyXSxbMSwxMiwiMV97Z2Z9IiwwLHsic2hvcnRlbiI6eyJzb3VyY2UiOjIwLCJ0YXJnZXQiOjIwfX1dLFs2LDE1LCIiLDAseyJzaG9ydGVuIjp7InNvdXJjZSI6MjAsInRhcmdldCI6MjB9LCJzdHlsZSI6eyJoZWFkIjp7Im5hbWUiOiJub25lIn19fV1d
		\[\begin{tikzcd}[ampersand replacement=\&]
			\& B \&\&\&\& {\st B} \& \\
			\& c \&\&\&\& c \\
			b \&\& d \& \rightsquigarrow \& b \&\& d
			\arrow["f"', from=1-2, to=1-6]
			\arrow["g", from=2-2, to=3-3]
			\arrow["g", from=2-6, to=3-7]
			\arrow["f", from=3-1, to=2-2]
			\arrow[""{name=0, anchor=center, inner sep=0}, "gf"', from=3-1, to=3-3]
			\arrow["f", from=3-5, to=2-6]
			\arrow[""{name=1, anchor=center, inner sep=0}, "gf"', from=3-5, to=3-7]
			\arrow[between={0.2}{0.8}, equals, nfold, from=2-2, to=0]
			\arrow["{1_{gf}}", between={0.2}{0.8}, Rightarrow, nfold, from=2-6, to=1]
		\end{tikzcd}\cdpunct[8pt]{,}\]
		\item[Identity constraint -] the $2$-cell $\varepsilon_{b} \Rightarrow 1_b$ is induced by the unique identity cell $1_{1_b} \colon 1_b \Rightarrow 1_b$,
		% https://q.uiver.app/#q=WzAsNyxbNCwyLCJiIl0sWzYsMiwiYiJdLFs1LDAsIlxcc3QgQiJdLFsxLDAsIkIiXSxbMCwyLCJiIl0sWzIsMiwiYiJdLFszLDIsIlxccmlnaHRzcXVpZ2Fycm93Il0sWzMsMiwiZiIsMl0sWzAsMSwiMV9iIiwyLHsiY3VydmUiOjJ9XSxbNCw1LCIxX2IiLDJdLFswLDEsIlxcdmFyZXBzaWxvbl97Yn0iLDAseyJjdXJ2ZSI6LTJ9XSxbMTAsOCwiMV97MV9ifSIsMCx7InNob3J0ZW4iOnsic291cmNlIjoyMCwidGFyZ2V0IjoyMH19XV0=
		\[\begin{tikzcd}[ampersand replacement=\&]
			\& B \&\&\&\& {\st B} \& \\
			\\
			b \&\& b \& \rightsquigarrow \& b \&\& b
			\arrow["f"', from=1-2, to=1-6]
			\arrow["{1_b}"', from=3-1, to=3-3]
			\arrow[""{name=0, anchor=center, inner sep=0}, "{1_b}"', curve={height=12pt}, from=3-5, to=3-7]
			\arrow[""{name=1, anchor=center, inner sep=0}, "{\varepsilon_{b}}", curve={height=-12pt}, from=3-5, to=3-7]
			\arrow["{1_{1_b}}", between={0.2}{0.8}, Rightarrow, nfold, from=1, to=0]
		\end{tikzcd}\cdpunct[8pt]{.}\]
	\end{description}
\end{definition}

It is highly important to note that the constraints cells of $s$ are non-trivial, even if both of them are represented by the identity $2$-cells under $e$. Furthermore, the equality $e \circ s = 1_{B}$ is meant as equality of pseudofunctors, which includes their coherence constraints. We will not verify the pseudo-functor axioms for both $e$ and $s$, as they are routine technical checks. 

\begin{comment}
	However, we do wish to remark that the coherence data of both $s$ and $e$ may have multiple presentations. For instance, the compositor of $e$ may also be presented by the following cell in $\st B$:
	% https://q.uiver.app/#q=WzAsOSxbMCwzLCJiIl0sWzIsMSwiYyJdLFsyLDAsIlxcc3QgQiJdLFs4LDAsIkIiXSxbNiwzLCJiIl0sWzgsMSwiYyJdLFsxMCwzLCJkIl0sWzUsMywiXFxtYXBzdG8iXSxbNCwzLCJkIl0sWzIsMywiZSJdLFs0LDUsImUoXFx2ZWN7Zn0pIl0sWzUsNiwiZShcXHZlY3tnfSkiXSxbNCw2LCJlKFxcdmVje2d9LFxcdmVje2Z9KSIsMl0sWzAsMSwiZShcXHZlY3tmfSkiXSxbMCw4LCJcXHZlY3tnfSxcXHZlY3tmfSIsMix7ImN1cnZlIjo1fV0sWzEsOCwiZShcXHZlY3tnfSkiXSxbMCw4LCJlKFxcdmVje2d9LFxcdmVje2Z9KSIsMSx7ImN1cnZlIjotMn1dLFs1LDEyLCIiLDAseyJzaG9ydGVuIjp7InNvdXJjZSI6MjAsInRhcmdldCI6MjB9fV0sWzEsMTYsIiIsMCx7InNob3J0ZW4iOnsic291cmNlIjoyMCwidGFyZ2V0IjoyMH19XSxbMTYsMTQsIjFfe2UoXFx2ZWN7Z30sXFx2ZWN7Zn0pfSIsMCx7InNob3J0ZW4iOnsic291cmNlIjoyMCwidGFyZ2V0IjoyMH19XV0=
	\[\begin{tikzcd}[ampersand replacement=\&]
		\&\& {\st B} \&\&\&\&\&\& B \&\& \\
		\&\& c \&\&\&\&\&\& c \\
		\\
		b \&\&\&\& d \& \mapsto \& b \&\&\&\& d
		\arrow["e", from=1-3, to=1-9]
		\arrow["{e(\vec{g})}", from=2-3, to=4-5]
		\arrow["{e(\vec{g})}", from=2-9, to=4-11]
		\arrow["{e(\vec{f})}", from=4-1, to=2-3]
		\arrow[""{name=0, anchor=center, inner sep=0}, "{\vec{g},\vec{f}}"', curve={height=30pt}, from=4-1, to=4-5]
		\arrow[""{name=1, anchor=center, inner sep=0}, "{e(\vec{g},\vec{f})}"{description}, curve={height=-12pt}, from=4-1, to=4-5]
		\arrow["{e(\vec{f})}", from=4-7, to=2-9]
		\arrow[""{name=2, anchor=center, inner sep=0}, "{e(\vec{g},\vec{f})}"', from=4-7, to=4-11]
		\arrow[between={0.2}{0.8}, Rightarrow, nfold, from=2-3, to=1]
		\arrow[between={0.2}{0.8}, Rightarrow, nfold, from=2-9, to=2]
		\arrow["{1_{e(\vec{g},\vec{f})}}", between={0.2}{0.8}, Rightarrow, nfold, from=1, to=0]
	\end{tikzcd}\cdpunct[8pt]{.}\]
	Ultimately, this is just a result of taking vertical composite with a $2$-cell whose image under $e$ is the identity. A similar result holds for the identity coherence cell of $e$ as well. In general, any composite of coherence data of $e$ may have multiple presentations.
\end{comment}

\subsection{\tpdf{Mapping of $\BiCat(A,B)$}}

Our goal now is determine the action of $\st$ on the bicategory of pseudofunctors, pseudotransformations, and modification. To this end, the following lemma shall be essential.

\begin{lemma}
	The dotted composite $\Phi$ defined in the diagram
	% https://q.uiver.app/#q=WzAsNSxbMCwwLCJcXEdyYXkoXFxzdCBBLEIpIl0sWzIsMCwiXFxCaUNhdChpKFxcc3QgQSksaUIpIl0sWzMsMCwiXFxCaUNhdChpKFxcc3QgQSksaUIpIFxcdGltZXMgKiJdLFszLDIsIlxcQmlDYXQoaShcXHN0IEEpLGlCKSBcXHRpbWVzIFxcQmlDYXQoQSxpKFxcc3QgQSkpIl0sWzIsMiwiXFxCaUNhdChBLGlCKSJdLFswLDEsImkiXSxbMSwyLCJcXGNvbmciXSxbMiwzLCIxIFxcdGltZXMgc197QX0iXSxbMyw0LCJcXGNpcmMiXSxbMCw0LCJcXFBoaSIsMix7InN0eWxlIjp7ImJvZHkiOnsibmFtZSI6ImRhc2hlZCJ9fX1dLFsxLDQsIihzX3tBfSleeyp9Il1d
	\[\begin{tikzcd}[ampersand replacement=\&,cramped]
		{\Gray(\st A,B)} \&\& {\BiCat(i(\st A),iB)} \& {\BiCat(i(\st A),iB) \times *} \\
		\\
		\&\& {\BiCat(A,iB)} \& {\BiCat(i(\st A),iB) \times \BiCat(A,i(\st A))}
		\arrow["i", from=1-1, to=1-3]
		\arrow["\Phi"', dashed, from=1-1, to=3-3]
		\arrow["\cong", from=1-3, to=1-4]
		\arrow["{(s_{A})^{*}}", from=1-3, to=3-3]
		\arrow["{1 \times s_{A}}", from=1-4, to=3-4]
		\arrow["\circ", from=3-4, to=3-3]
	\end{tikzcd}\]
	is a strict $2$-functor of strict $2$-categories. Here, $s_{A}$ is the constant pseudofunctor out of the terminal bicategory which picks $s_{A} \in \BiCat(A,i(\st A))$, and the functor $\circ$ is the compositor of $\BiCat$.
\end{lemma}

\begin{proof}
	We first note that $\BiCat(A,Y)$ is a strict $2$-category, whenever $Y$ is a strict $2$-category. Thus, all of the bicategories involved in defining the above composite are all strict $2$-categories. Next, we verify that $\Phi$ is a strict $2$-functor of strict $2$-categories. Note that $\Phi(\gamma \colon F \Rightarrow  G)$ is given by $\gamma \star 1_{S_{A}}$. Using the standard definition of the compositor in $\BiCat$, we can represent this using the following transformation cell in $B$:
	% https://q.uiver.app/#q=WzAsMTMsWzIsMCwiRmEiXSxbNCwwLCJGYiJdLFsyLDIsIkdhIl0sWzQsMiwiR2IiXSxbNSwyLCJcXGVxdWl2Il0sWzYsMSwiRmEiXSxbNiwzLCJHYSJdLFs4LDEsIkZiIl0sWzgsMywiR2IiXSxbMiw0LCJHYSJdLFs0LDQsIkdiIl0sWzAsMiwiXFxQaGkoXFxnYW1tYSkiXSxbMSwyLCJcXGVxdWl2Il0sWzAsMSwiRmYiXSxbMCwyLCJcXGdhbW1hX3thfSIsMl0sWzEsMywiXFxnYW1tYV97Yn0iXSxbMiwzLCJHZiIsMV0sWzIsMSwiXFxnYW1tYV97Zn0iLDEseyJzaG9ydGVuIjp7InNvdXJjZSI6MjAsInRhcmdldCI6MjB9LCJsZXZlbCI6Mn1dLFs2LDgsIkdmIiwyXSxbNSw2LCJcXGdhbW1hX2EiLDJdLFs1LDcsIkZmIl0sWzcsOCwiXFxnYW1tYV9iIl0sWzYsNywiXFxnYW1tYV9mIiwxLHsic2hvcnRlbiI6eyJzb3VyY2UiOjIwLCJ0YXJnZXQiOjIwfSwibGV2ZWwiOjJ9XSxbMiw5LCIxX3tHYX0iLDJdLFszLDEwLCIxX3tHYX0iXSxbOSwxMCwiR2YiLDJdLFs5LDMsIjFfe0dmfSIsMSx7InNob3J0ZW4iOnsic291cmNlIjoyMCwidGFyZ2V0IjoyMH0sImxldmVsIjoyfV1d
	\[\begin{tikzcd}[ampersand replacement=\&]
		\&\& Fa \&\& Fb \&\&\&\& \\
		\&\&\&\&\&\& Fa \&\& Fb \\
		{\Phi(\gamma)} \& \equiv \& Ga \&\& Gb \& \equiv \\
		\&\&\&\&\&\& Ga \&\& Gb \\
		\&\& Ga \&\& Gb
		\arrow["Ff", from=1-3, to=1-5]
		\arrow["{\gamma_{a}}"', from=1-3, to=3-3]
		\arrow["{\gamma_{b}}", from=1-5, to=3-5]
		\arrow["Ff", from=2-7, to=2-9]
		\arrow["{\gamma_a}"', from=2-7, to=4-7]
		\arrow["{\gamma_b}", from=2-9, to=4-9]
		\arrow["{\gamma_{f}}"{description}, between={0.2}{0.8}, Rightarrow, nfold, from=3-3, to=1-5]
		\arrow["Gf"{description}, from=3-3, to=3-5]
		\arrow["{1_{Ga}}"', from=3-3, to=5-3]
		\arrow["{1_{Ga}}", from=3-5, to=5-5]
		\arrow["{\gamma_f}"{description}, between={0.2}{0.8}, Rightarrow, nfold, from=4-7, to=2-9]
		\arrow["Gf"', from=4-7, to=4-9]
		\arrow["{1_{Gf}}"{description}, between={0.2}{0.8}, Rightarrow, nfold, from=5-3, to=3-5]
		\arrow["Gf"', from=5-3, to=5-5]
	\end{tikzcd}\cdpunct[8pt]{.}\]
	Crucially, we emphasize that this equality is possible only because both $F,G \colon \st A \rightarrow B$ are strict $2$-functors of strict $2$-categories. With this description, it is now straightforward to verify both $(\gamma_2\gamma_1) \star 1_{S_{A}}=(\gamma_2 \star 1_{S_{A}})(\gamma_1 \star 1_{S_{A}})$ and $1_F \star 1_{S_{A}} = 1_{F \star S_{A}}$. Thus, the pseudofunctor $\Phi$ is indeed strict $2$-functor.
	
\end{proof}

It turns out that $\Phi$ is actually a strict $2$-isomorphism of $2$-categories. For this, we only need to show that it is a bijection on the underlying $2$-graphs. Thus, we state the following lemma.

\begin{lemma} \leavevmode
	\begin{enumerate}
		\item For every pseudofunctor $U\colon A \rightarrow Y$ with $Y \in \Gray$, there exists a unique strict functor $\tilde U \colon \st A \rightarrow Y$ such that $\tilde{U} \circ s_{A} = U$.
		\item For every pseudotransformation $\gamma\colon U \Rightarrow V \colon A \rightarrow Y$ with $Y \in \Gray$, there exists a unique pseudotransformation $\tilde \gamma$ such that $\tilde{\gamma} \star 1_{s_{A}} = \rho$.
		\item For every modification $\theta \colon \gamma \Rrightarrow \delta \colon U \Rightarrow V  \colon A \rightarrow Y$ with $Y \in \Gray$, there exists a unique modification $\tilde{\theta}$ such that $\tilde{\theta} \cdot 1_{1_{s_{A}}} = \omega$.
	\end{enumerate}
\end{lemma}

\begin{proof} 
	The only non-trivial part of the proof is defining each of the unique lifts, the rest is a routine.
	\begin{enumerate}
		\item Define $\tilde{U}$ by the following data:
		\begin{description}
			\item[Objects -] for any $a \in \st A$, $U$ maps $a \mapsto Ua \in Y$,
			\item[$1$-cells -] the non-identity $1$-cells $\vec{f}$ are mapped as
			% https://q.uiver.app/#q=WzAsNyxbMSwwLCJhXzEiXSxbMiwwLCJcXGNkb3RzIl0sWzAsMCwiYSJdLFszLDAsIlxcbWFwc3RvXntcXHRpbGRle1V9fSJdLFs0LDAsIlVhIl0sWzUsMCwiVWFfMSJdLFs2LDAsIlxcY2RvdHMiXSxbMCwxLCJmXzIiXSxbMiwwLCJmXzEiXSxbNCw1LCJVZl8xIl0sWzUsNiwiVWZfMiJdXQ==
			\[\begin{tikzcd}[ampersand replacement=\&,cramped]
				a \& {a_1} \& \cdots \& {\mapsto^{\tilde{U}}} \& Ua \& {Ua_1} \& \cdots
				\arrow["{f_1}", from=1-1, to=1-2]
				\arrow["{f_2}", from=1-2, to=1-3]
				\arrow["{Uf_1}", from=1-5, to=1-6]
				\arrow["{Uf_2}", from=1-6, to=1-7]
			\end{tikzcd}\cdpunct[8pt]{,}\]
			and the identity $1$-cell is mapped to the identity $1$-cell,
			\item[$2$-cells -] the action on $2$-cells is given as
			% https://q.uiver.app/#q=WzAsMzMsWzEsNCwiYV8xIl0sWzIsNCwiXFxjZG90cyJdLFswLDQsImEiXSxbNSw0LCJVYSJdLFs2LDQsIlVhXzEiXSxbNyw0LCJcXGNkb3RzIl0sWzAsOCwiYSJdLFsxLDgsImEnXzEiXSxbMiw4LCJcXGNkb3RzIl0sWzUsOCwiVWEiXSxbNiw4LCJVYSdfMSJdLFs3LDgsIlxcY2RvdHMiXSxbMCw1LCJhIl0sWzMsNSwiYiJdLFszLDQsImIiXSxbMyw4LCJiIl0sWzAsNywiYSJdLFszLDcsImIiXSxbMCwwLCJhIl0sWzEsMCwiYV8xIl0sWzIsMCwiXFxjZG90cyJdLFszLDAsImIiXSxbMCwyLCJhIl0sWzEsMiwiYSdfMSJdLFsyLDIsIlxcY2RvdHMiXSxbMywyLCJiIl0sWzUsNSwiVWEiXSxbOCw0LCJVYiJdLFs4LDgsIlViIl0sWzUsNywiVWEiXSxbOCw3LCJVYiJdLFs4LDUsIlViIl0sWzQsNiwiXFxtYXBzdG9ee1xcdGlsZGV7VX19Il0sWzAsMSwiZl8yIl0sWzIsMCwiZl8xIl0sWzMsNCwiVWZfMSJdLFs0LDUsIlVmXzIiXSxbNyw4LCJnXzIiLDJdLFs2LDcsImdfMSIsMl0sWzksMTAsIlVnXzEiLDJdLFsxMCwxMSwiVWdfMiIsMl0sWzgsMTVdLFsxLDE0XSxbMTIsMTMsImUoXFx2ZWN7Zn0pIiwxXSxbMTYsMTcsImUoXFx2ZWN7Z30pIiwxXSxbMjIsMjMsImdfMSIsMl0sWzIzLDI0LCJnXzIiLDJdLFsyNCwyNV0sWzE4LDE5LCJmXzEiXSxbMTksMjAsImZfMiJdLFsyMCwyMV0sWzUsMjddLFsxMSwyOF0sWzI2LDMxLCJVKGUoXFx2ZWN7Zn0pKSIsMV0sWzI5LDMwLCJVKGUoXFx2ZWN7Z30pKSIsMV0sWzMzLDQzLCJcXG1hdGhybXtDb2h9X3tTX3tBfX0iLDAseyJzaG9ydGVuIjp7InNvdXJjZSI6MjAsInRhcmdldCI6MjB9fV0sWzQ0LDM3LCJcXG1hdGhybXtDb2h9X3tTX3tBfX1eey0xfSIsMCx7InNob3J0ZW4iOnsic291cmNlIjoyMCwidGFyZ2V0IjoyMH19XSxbNDksNDYsIihcXGFscGhhIDogZShcXHZlY3tmfSkgXFxSaWdodGFycm93IGUoXFx2ZWN7Z30pKSIsMCx7InNob3J0ZW4iOnsic291cmNlIjoyMCwidGFyZ2V0IjoyMH19XSxbNDMsNDQsIlxcYWxwaGEiLDEseyJzaG9ydGVuIjp7InNvdXJjZSI6MjAsInRhcmdldCI6MjB9fV0sWzQ2LDMzLCJcXGVxdWl2IiwzLHsic2hvcnRlbiI6eyJzb3VyY2UiOjIwLCJ0YXJnZXQiOjIwfSwic3R5bGUiOnsiYm9keSI6eyJuYW1lIjoibm9uZSJ9LCJoZWFkIjp7Im5hbWUiOiJub25lIn19fV0sWzUzLDU0LCJVKFxcYWxwaGEpIiwwLHsic2hvcnRlbiI6eyJzb3VyY2UiOjIwLCJ0YXJnZXQiOjIwfX1dLFs1NCw0MCwiXFxtYXRocm17Q29ofV9VXnstMX0iLDAseyJzaG9ydGVuIjp7InNvdXJjZSI6MjAsInRhcmdldCI6MjB9fV0sWzM2LDUzLCJcXG1hdGhybXtDb2h9X1UiLDAseyJzaG9ydGVuIjp7InNvdXJjZSI6MjAsInRhcmdldCI6MjB9fV1d
			\[\begin{tikzcd}[ampersand replacement=\&,cramped]
				a \& {a_1} \& \cdots \& b \&\&\&\&\& \\
				\\
				a \& {a'_1} \& \cdots \& b \\
				\\
				a \& {a_1} \& \cdots \& b \&\& Ua \& {Ua_1} \& \cdots \& Ub \\
				a \&\&\& b \&\& Ua \&\&\& Ub \\
				\&\&\&\& {\mapsto^{\tilde{U}}} \\
				a \&\&\& b \&\& Ua \&\&\& Ub \\
				a \& {a'_1} \& \cdots \& b \&\& Ua \& {Ua'_1} \& \cdots \& Ub
				\arrow["{f_1}", from=1-1, to=1-2]
				\arrow[""{name=0, anchor=center, inner sep=0}, "{f_2}", from=1-2, to=1-3]
				\arrow[from=1-3, to=1-4]
				\arrow["{g_1}"', from=3-1, to=3-2]
				\arrow[""{name=1, anchor=center, inner sep=0}, "{g_2}"', from=3-2, to=3-3]
				\arrow[from=3-3, to=3-4]
				\arrow["{f_1}", from=5-1, to=5-2]
				\arrow[""{name=2, anchor=center, inner sep=0}, "{f_2}", from=5-2, to=5-3]
				\arrow[from=5-3, to=5-4]
				\arrow["{Uf_1}", from=5-6, to=5-7]
				\arrow[""{name=3, anchor=center, inner sep=0}, "{Uf_2}", from=5-7, to=5-8]
				\arrow[from=5-8, to=5-9]
				\arrow[""{name=4, anchor=center, inner sep=0}, "{e(\vec{f})}"{description}, from=6-1, to=6-4]
				\arrow[""{name=5, anchor=center, inner sep=0}, "{U(e(\vec{f}))}"{description}, from=6-6, to=6-9]
				\arrow[""{name=6, anchor=center, inner sep=0}, "{e(\vec{g})}"{description}, from=8-1, to=8-4]
				\arrow[""{name=7, anchor=center, inner sep=0}, "{U(e(\vec{g}))}"{description}, from=8-6, to=8-9]
				\arrow["{g_1}"', from=9-1, to=9-2]
				\arrow[""{name=8, anchor=center, inner sep=0}, "{g_2}"', from=9-2, to=9-3]
				\arrow[from=9-3, to=9-4]
				\arrow["{Ug_1}"', from=9-6, to=9-7]
				\arrow[""{name=9, anchor=center, inner sep=0}, "{Ug_2}"', from=9-7, to=9-8]
				\arrow[from=9-8, to=9-9]
				\arrow["{(\alpha : e(\vec{f}) \Rightarrow e(\vec{g}))}", between={0.2}{0.8}, Rightarrow, nfold, from=0, to=1]
				\arrow["\equiv"{marking, allow upside down}, draw=none, from=1, to=2]
				\arrow["{\mathrm{Coh}_{S_{A}}}", between={0.2}{0.8}, Rightarrow, nfold, from=2, to=4]
				\arrow["{\mathrm{Coh}_U}", between={0.2}{0.8}, Rightarrow, nfold, from=3, to=5]
				\arrow["\alpha"{description}, between={0.2}{0.8}, Rightarrow, nfold, from=4, to=6]
				\arrow["{U(\alpha)}", between={0.2}{0.8}, Rightarrow, nfold, from=5, to=7]
				\arrow["{\mathrm{Coh}_{S_{A}}^{-1}}", between={0.2}{0.8}, Rightarrow, nfold, from=6, to=8]
				\arrow["{\mathrm{Coh}_U^{-1}}", between={0.2}{0.8}, Rightarrow, nfold, from=7, to=9]
			\end{tikzcd}\cdpunct[8pt]{,}\]
			where $\mathrm{Coh}_U$ and $\mathrm{Coh}_{S_{A}}$ represent the unique coherence constraint cells.
		\end{description}
		
		By choosing $1_{\vec{f}} \colon \vec{f} \Rightarrow e(\vec{f})$, the definition readily implies that $\tilde{U}(\mathrm{Coh}_{s_{A}}) = \mathrm{Coh}_U$. Thus, $\tilde{U} \circ s_{A} = U$ holds as desired, including the equality of coherence data.
		
		\item Define $\tilde\gamma$ by the following data:
		\begin{description}
			\item[Component $1$-cell] for each $a \in \st A$, the component $(\st \gamma)_{a}$ is defined as $\gamma_{a}$ itself,
			\item[Component $2$-cell] for each string $\vec{f} \in \st A$, the component $(\st \gamma)_{\vec{f}}$ is defined by the following composite $2$-cell in $Y$
			% https://q.uiver.app/#q=WzAsMTAsWzAsMCwiVWEiXSxbMCwyLCJWYSJdLFs0LDAsIlViIl0sWzQsMiwiVmIiXSxbMSwyLCJWYV8xIl0sWzIsMiwiVmFfMiJdLFszLDIsIlxcY2RvdHMiXSxbMSwwLCJVYV8xIl0sWzIsMCwiVWFfMiJdLFszLDAsIlxcY2RvdHMiXSxbMSw0LCJWZl8xIiwyXSxbNCw1LCJWZl8yIiwyXSxbNSw2LCJWZl8zIiwyXSxbNiwzXSxbMCw3LCJVZl8xIl0sWzcsOCwiVWZfMiJdLFs4LDksIlVmXzMiXSxbOSwyXSxbMCwxLCJcXGdhbW1hX2EiLDJdLFs3LDQsIlxcZ2FtbWFfe2FfMX0iLDFdLFs4LDUsIlxcZ2FtbWFfe2FfMn0iLDFdLFsyLDMsIlxcZ2FtbWFfYiJdLFsxLDcsIlxcZ2FtbWFfe2ZfMX0iLDEseyJzaG9ydGVuIjp7InNvdXJjZSI6MjAsInRhcmdldCI6MjB9LCJsZXZlbCI6Mn1dLFs0LDgsIlxcZ2FtbWFfe2ZfMn0iLDEseyJzaG9ydGVuIjp7InNvdXJjZSI6MjAsInRhcmdldCI6MjB9LCJsZXZlbCI6Mn1dXQ==
			\[\begin{tikzcd}[ampersand replacement=\&,cramped]
				Ua \& {Ua_1} \& {Ua_2} \& \cdots \& Ub \\
				\\
				Va \& {Va_1} \& {Va_2} \& \cdots \& Vb
				\arrow["{Uf_1}", from=1-1, to=1-2]
				\arrow["{\gamma_a}"', from=1-1, to=3-1]
				\arrow["{Uf_2}", from=1-2, to=1-3]
				\arrow["{\gamma_{a_1}}"{description}, from=1-2, to=3-2]
				\arrow["{Uf_3}", from=1-3, to=1-4]
				\arrow["{\gamma_{a_2}}"{description}, from=1-3, to=3-3]
				\arrow[from=1-4, to=1-5]
				\arrow["{\gamma_b}", from=1-5, to=3-5]
				\arrow["{\gamma_{f_1}}"{description}, between={0.2}{0.8}, Rightarrow, nfold, from=3-1, to=1-2]
				\arrow["{Vf_1}"', from=3-1, to=3-2]
				\arrow["{\gamma_{f_2}}"{description}, between={0.2}{0.8}, Rightarrow, nfold, from=3-2, to=1-3]
				\arrow["{Vf_2}"', from=3-2, to=3-3]
				\arrow["{Vf_3}"', from=3-3, to=3-4]
				\arrow[from=3-4, to=3-5]
			\end{tikzcd}\cdpunct[8pt]{.}\]
		\end{description}
		
		\item Define $\tilde \theta \colon \tilde{\gamma} \Rrightarrow \tilde{\delta}$ by setting $(\tilde \theta)_{a} = \theta_a \colon \gamma_a \Rightarrow \delta_a$
	\end{enumerate}
\end{proof}

As a consequence, we obtain the desired proposition.

\begin{proposition}
		The dotted composite $\Phi$ defined in the diagram
		% https://q.uiver.app/#q=WzAsNSxbMCwwLCJcXEdyYXkoXFxzdCBBLEIpIl0sWzIsMCwiXFxCaUNhdChpKFxcc3QgQSksaUIpIl0sWzMsMCwiXFxCaUNhdChpKFxcc3QgQSksaUIpIFxcdGltZXMgKiJdLFszLDIsIlxcQmlDYXQoaShcXHN0IEEpLGlCKSBcXHRpbWVzIFxcQmlDYXQoQSxpKFxcc3QgQSkpIl0sWzIsMiwiXFxCaUNhdChBLGlCKSJdLFswLDEsImkiXSxbMSwyLCJcXGNvbmciXSxbMiwzLCIxIFxcdGltZXMgc197QX0iXSxbMyw0LCJcXGNpcmMiXSxbMCw0LCJcXFBoaSIsMix7InN0eWxlIjp7ImJvZHkiOnsibmFtZSI6ImRhc2hlZCJ9fX1dLFsxLDQsIihzX3tBfSleeyp9Il1d
		\[\begin{tikzcd}[ampersand replacement=\&,cramped]
			{\Gray(\st A,B)} \&\& {\BiCat(i(\st A),iB)} \& {\BiCat(i(\st A),iB) \times *} \\
			\\
			\&\& {\BiCat(A,iB)} \& {\BiCat(i(\st A),iB) \times \BiCat(A,i(\st A))}
			\arrow["i", from=1-1, to=1-3]
			\arrow["\Phi"', dashed, from=1-1, to=3-3]
			\arrow["\cong", from=1-3, to=1-4]
			\arrow["{(s_{A})^{*}}", from=1-3, to=3-3]
			\arrow["{1 \times s_{A}}", from=1-4, to=3-4]
			\arrow["\circ", from=3-4, to=3-3]
		\end{tikzcd}\]
		is a strict $2$-isomorphism of strict $2$-categories.
\end{proposition}

\begin{remark}
	One may visualize this universal property using the diagram 
	% https://q.uiver.app/#q=WzAsMyxbMCwwLCJpIFxcY2lyYyBcXHN0IEEiXSxbOCwwLCJpWSJdLFswLDYsIkEiXSxbMCwxLCJcXHRpbGRle1V9IiwwLHsiY3VydmUiOi00fV0sWzIsMCwic197QX0iXSxbMiwxLCJVIiwwLHsiY3VydmUiOi00fV0sWzIsMSwiViIsMix7ImN1cnZlIjo0fV0sWzAsMSwiXFx0aWxkZXtWfSIsMix7ImN1cnZlIjo0fV0sWzMsNywiXFx0aWxkZXtcXGdhbW1hfSIsMix7ImN1cnZlIjo0LCJzaG9ydGVuIjp7InNvdXJjZSI6MjAsInRhcmdldCI6MjB9fV0sWzUsNiwiXFxnYW1tYSIsMix7ImN1cnZlIjo0LCJzaG9ydGVuIjp7InNvdXJjZSI6MjAsInRhcmdldCI6MjB9fV0sWzUsNiwiXFxkZWx0YSIsMCx7ImN1cnZlIjotNCwic2hvcnRlbiI6eyJzb3VyY2UiOjIwLCJ0YXJnZXQiOjIwfX1dLFszLDcsIlxcdGlsZGV7XFxkZWx0YX0iLDAseyJjdXJ2ZSI6LTQsInNob3J0ZW4iOnsic291cmNlIjoyMCwidGFyZ2V0IjoyMH19XSxbOSwxMCwiXFx0aGV0YSIsMCx7InNob3J0ZW4iOnsic291cmNlIjoyMCwidGFyZ2V0IjoyMH19XSxbOCwxMSwiXFx0aWxkZXtcXHRoZXRhfSIsMCx7InNob3J0ZW4iOnsic291cmNlIjoyMCwidGFyZ2V0IjoyMH19XV0=
	\[\begin{tikzcd}[ampersand replacement=\&]
		{i \circ \st A} \&\&\&\&\&\&\&\& iY \\
		\\
		\\
		\\
		\\
		\\
		A
		\arrow[""{name=0, anchor=center, inner sep=0}, "{\tilde{U}}", curve={height=-24pt}, from=1-1, to=1-9]
		\arrow[""{name=1, anchor=center, inner sep=0}, "{\tilde{V}}"', curve={height=24pt}, from=1-1, to=1-9]
		\arrow["{s_{A}}", from=7-1, to=1-1]
		\arrow[""{name=2, anchor=center, inner sep=0}, "U", curve={height=-24pt}, from=7-1, to=1-9]
		\arrow[""{name=3, anchor=center, inner sep=0}, "V"', curve={height=24pt}, from=7-1, to=1-9]
		\arrow[""{name=4, anchor=center, inner sep=0}, "{\tilde{\gamma}}"', curve={height=24pt}, between={0.2}{0.8}, Rightarrow, nfold, from=0, to=1]
		\arrow[""{name=5, anchor=center, inner sep=0}, "{\tilde{\delta}}", curve={height=-24pt}, between={0.2}{0.8}, Rightarrow, nfold, from=0, to=1]
		\arrow[""{name=6, anchor=center, inner sep=0}, "\gamma"', curve={height=24pt}, between={0.2}{0.8}, Rightarrow, nfold, from=2, to=3]
		\arrow[""{name=7, anchor=center, inner sep=0}, "\delta", curve={height=-24pt}, between={0.2}{0.8}, Rightarrow, nfold, from=2, to=3]
		\arrow["{\tilde{\theta}}", between={0.2}{0.8}, Rightarrow, scaling nfold=3, from=4, to=5]
		\arrow["\theta", between={0.2}{0.8}, Rightarrow, scaling nfold=3, from=6, to=7]
	\end{tikzcd}\cdpunct[8pt]{.}\]
\end{remark}

As with any ordinary adjunction, we can utilize this universal property to extend $\st$ to a trifunctor. We first extend it on $\BiCat(A,B)$.

\begin{definition}
	Define the post composition pseudofunctor $(s_{B})_{*}$, corresponding to $s_{B} \colon B \rightarrow \st B$, as the composite pseudofunctor
	% https://q.uiver.app/#q=WzAsNCxbMCwwLCJcXEJpQ2F0KEEsQikiXSxbMiwwLCIqIFxcdGltZXMgXFxCaUNhdChBLEIpICJdLFsyLDIsIlxcQmlDYXQoQixpXFxjaXJjXFxzdCBCKSBcXHRpbWVzIFxcQmlDYXQoQSxCKSJdLFswLDIsIlxcQmlDYXQoQSxpIFxcY2lyYyBcXHN0IEIpIl0sWzAsMSwiXFxjb25nIl0sWzEsMiwic197Qn0gXFx0aW1lcyAxIl0sWzIsMywiXFxjaXJjIl0sWzAsMywiKHNfe0J9KV97Kn0iLDJdXQ==
	\[\begin{tikzcd}[ampersand replacement=\&]
		{\BiCat(A,B)} \&\& {* \times \BiCat(A,B) } \\
		\\
		{\BiCat(A,i \circ \st B)} \&\& {\BiCat(B,i\circ\st B) \times \BiCat(A,B)}
		\arrow["\cong", from=1-1, to=1-3]
		\arrow["{(s_{B})_{*}}"', from=1-1, to=3-1]
		\arrow["{s_{B} \times 1}", from=1-3, to=3-3]
		\arrow["\circ", from=3-3, to=3-1]
	\end{tikzcd}\cdpunct[8pt]{.}\]
\end{definition}

Note that $(s_{B})_{*}$ is not a strict $2$-functor. The compositor is a modification given by unique coherence cell $(1_{s_{B}} \star \gamma_2) (1_{s_{B}} \star \gamma_1) \Rrightarrow (1_{s_{B}} \star \gamma_2\gamma_1)$. Similarly, the identity coherence modification is given by the unique coherence cell $1_{s_{B} \circ F} \Rrightarrow 1_{s_{B}} \star 1_F$.

\begin{definition}
	Define $\st_{A,B} \colon \BiCat(A,B) \rightarrow \Gray(\st A, \st B)$ as the composite pseudofunctor
	% https://q.uiver.app/#q=WzAsMyxbMCwwLCJcXEJpQ2F0KEEsQikiXSxbMiwwLCJcXEJpQ2F0KEEsaSBcXGNpcmMgXFxzdCBCKSJdLFs0LDAsIlxcR3JheShcXHN0IEEsIFxcc3QgQikiXSxbMCwxLCIoc197Qn0pX3sqfSJdLFsxLDIsIlxcUGhpXnstMX0iXV0=
	\[\begin{tikzcd}[ampersand replacement=\&]
		{\BiCat(A,B)} \&\& {\BiCat(A,i \circ \st B)} \&\& {\Gray(\st A, \st B)}
		\arrow["{(s_{B})_{*}}", from=1-1, to=1-3]
		\arrow["{\Phi^{-1}}", from=1-3, to=1-5]
	\end{tikzcd}\cdpunct[1pt]{.}\]
\end{definition}

With these definitions, one can readily check that the following diagram commutes strictly
% https://q.uiver.app/#q=WzAsNCxbMCwwLCJcXEJpQ2F0KEEsQikiXSxbMywwLCJcXEJpQ2F0KGlcXGNpcmMgXFxzdCBBLCBpIFxcY2lyYyBcXHN0IEIpIl0sWzAsMywiXFxCaUNhdChBLEIpIl0sWzMsMywiXFxCaUNhdChBLGkgXFxjaXJjIFxcc3QgQikiXSxbMCwxLCJpX3tcXHN0IEEsXFxzdCBCfSBcXGNpcmMgXFxzdF97QSxCfSJdLFswLDIsIjEiLDJdLFsyLDMsIihzX3tCfSlfeyp9IiwyXSxbMSwzLCIoc197QX0pXnsqfSJdXQ==
\[\begin{tikzcd}[ampersand replacement=\&]
	{\BiCat(A,B)} \&\&\& {\BiCat(i\circ \st A, i \circ \st B)} \\
	\\
	\\
	{\BiCat(A,B)} \&\&\& {\BiCat(A,i \circ \st B)}
	\arrow["{i_{\st A,\st B} \circ \st_{A,B}}", from=1-1, to=1-4]
	\arrow["1"', from=1-1, to=4-1]
	\arrow["{(s_{A})^{*}}", from=1-4, to=4-4]
	\arrow["{(s_{B})_{*}}"', from=4-1, to=4-4]
\end{tikzcd}\cdpunct[8pt]{.}\]

Finally, the coherence data of $\st_{A,B}$ is simply the image under $\Phi^{-1}$ of the coherence constraint cells of $(s_{B})_{*}$. Thus, we obtain the following concrete description of $\st_{A,B}$.

\begin{comment}
	There is another useful characterization of the lifted pseudonatural transformation $\tilde{\gamma}$ which follows directly from defining axioms of a pseudonatural transformation and induction.
	
	\begin{lemma}\label{lem:st2pstrans}
		For any string $\vec{f} \in \st A$, the component $2$-cell $(\st \gamma)_{\vec{f}}$ is equivalently the following composite $2$-cell in $Y$:
		% https://q.uiver.app/#q=WzAsMTAsWzAsMiwiVWEiXSxbMCw0LCJWYSJdLFs0LDIsIlViIl0sWzQsNCwiVmIiXSxbMSw2LCJWYV8xIl0sWzIsNiwiVmFfMiJdLFszLDYsIlxcY2RvdHMiXSxbMSwwLCJVYV97MX0iXSxbMiwwLCJVYV8yIl0sWzMsMCwiXFxjZG90cyJdLFsxLDQsIlZmXzEiLDJdLFs0LDUsIlZmXzIiLDJdLFs1LDYsIlZmXzMiLDJdLFs2LDNdLFswLDcsIlVmXzEiXSxbNyw4LCJVZl8yIl0sWzgsOSwiVWZfMyJdLFs5LDJdLFswLDEsIlxcZ2FtbWFfYSIsMl0sWzIsMywiXFxnYW1tYV9iIl0sWzEsMywiVihlKFxcdmVje2Z9KSkiLDFdLFswLDIsIlUoZShcXHZlY3tmfSkpIiwxXSxbMSwyLCJcXGdhbW1hX3tlKFxcdmVje2Z9KX0iLDEseyJzaG9ydGVuIjp7InNvdXJjZSI6MjAsInRhcmdldCI6MjB9LCJsZXZlbCI6Mn1dLFs1LDIwLCJcXG1hdGhybXtDb2h9X1YiLDIseyJzaG9ydGVuIjp7InRhcmdldCI6MjB9fV0sWzIxLDgsIlxcbWF0aHJte0NvaH1fVV57LTF9IiwyLHsibGFiZWxfcG9zaXRpb24iOjYwLCJzaG9ydGVuIjp7InNvdXJjZSI6MjB9fV1d
		\[\begin{tikzcd}[ampersand replacement=\&]
			\& {Ua_{1}} \& {Ua_2} \& \cdots \& \\
			\\
			Ua \&\&\&\& Ub \\
			\\
			Va \&\&\&\& Vb \\
			\\
			\& {Va_1} \& {Va_2} \& \cdots
			\arrow["{Uf_2}", from=1-2, to=1-3]
			\arrow["{Uf_3}", from=1-3, to=1-4]
			\arrow[from=1-4, to=3-5]
			\arrow["{Uf_1}", from=3-1, to=1-2]
			\arrow[""{name=0, anchor=center, inner sep=0}, "{U(e(\vec{f}))}"{description}, from=3-1, to=3-5]
			\arrow["{\gamma_a}"', from=3-1, to=5-1]
			\arrow["{\gamma_b}", from=3-5, to=5-5]
			\arrow["{\gamma_{e(\vec{f})}}"{description}, between={0.2}{0.8}, Rightarrow, nfold, from=5-1, to=3-5]
			\arrow[""{name=1, anchor=center, inner sep=0}, "{V(e(\vec{f}))}"{description}, from=5-1, to=5-5]
			\arrow["{Vf_1}"', from=5-1, to=7-2]
			\arrow["{Vf_2}"', from=7-2, to=7-3]
			\arrow["{Vf_3}"', from=7-3, to=7-4]
			\arrow[from=7-4, to=5-5]
			\arrow["{\mathrm{Coh}_U^{-1}}"'{pos=0.6}, between={0.2}{1}, Rightarrow, nfold, from=0, to=1-3]
			\arrow["{\mathrm{Coh}_V}"', between={0}{0.8}, Rightarrow, nfold, from=7-3, to=1]
		\end{tikzcd}\]
	\end{lemma}
\end{comment}

\begin{lemma}
	The coherence data of $\st_{A,B} \colon \BiCat(A,B) \rightarrow \Gray(\st A,\st B)$ is given by
	\begin{description}
		\item[Composition constraint -] for every $\gamma_1 \colon F \Rightarrow G$ and $\gamma_2 \colon G \Rightarrow H$, $\st (\gamma_2)\st(\gamma_1) \Rrightarrow \st(\gamma_2\gamma_1)$ is given by a $2$-cell/modification in $\Gray(\st A, \st B)$, whose component at $a \in \st A$ is as depicted in:
		% https://q.uiver.app/#q=WzAsMyxbMSwwLCJGYSJdLFswLDIsIkdhIl0sWzEsNCwiSGEiXSxbMCwxLCIoXFxnYW1tYV8xKV9hIiwyXSxbMSwyLCIoXFxnYW1tYV8yKV9hIiwyXSxbMCwyLCIoXFxnYW1tYV8yKV9hIChcXGdhbW1hXzEpX2EiXSxbMSw1LCIiLDIseyJzaG9ydGVuIjp7InNvdXJjZSI6MjAsInRhcmdldCI6MjB9fV1d
		\[\begin{tikzcd}[ampersand replacement=\&]
			\& Fa \\
			\\
			Ga \\
			\\
			\& Ha
			\arrow["{(\gamma_1)_a}"', from=1-2, to=3-1]
			\arrow[""{name=0, anchor=center, inner sep=0}, "{(\gamma_2)_a (\gamma_1)_a}", from=1-2, to=5-2]
			\arrow["{(\gamma_2)_a}"', from=3-1, to=5-2]
			\arrow[between={0.2}{0.8}, Rightarrow, nfold, from=3-1, to=0]
		\end{tikzcd}\cdpunct[8pt]{,}\]
		\item[Identity constraint -]   for every $1_F\colon F \Rightarrow F \colon A \rightarrow B$, $1_{\st F} \Rrightarrow \st (1_F)$ is given by a $2$-cell/modification in $\Gray(\st A, \st B)$, whose component at $a \in \st A$ is as depicted in:
		% https://q.uiver.app/#q=WzAsMixbMCwwLCJGYSJdLFswLDIsIkZhIl0sWzAsMSwiXFx2YXJlcHNpbG9uX3tGYX0iLDIseyJjdXJ2ZSI6M31dLFswLDEsIjFfe0ZhfSIsMCx7ImN1cnZlIjotM31dLFsyLDMsIiIsMix7InNob3J0ZW4iOnsic291cmNlIjoyMCwidGFyZ2V0IjoyMH19XV0=
		\[\begin{tikzcd}[ampersand replacement=\&]
			Fa \\
			\\
			Fa
			\arrow[""{name=0, anchor=center, inner sep=0}, "{\varepsilon_{Fa}}"', curve={height=18pt}, from=1-1, to=3-1]
			\arrow[""{name=1, anchor=center, inner sep=0}, "{1_{Fa}}", curve={height=-18pt}, from=1-1, to=3-1]
			\arrow[between={0.2}{0.8}, Rightarrow, nfold, from=0, to=1]
		\end{tikzcd}\cdpunct[8pt]{.}\]
	\end{description}
\end{lemma}

\subsection{\tpdf{Coherence data of $\st$ as trifunctor}}
 
 To specify the functor $\st$ as a trifunctor, we also need to specify its identity and compositor adjoint equivalences. Note that we have not yet imposed any naturality condition on the isomorphisms $\Phi_{A,Y} \colon \Gray(\st A, Y) \rightarrow \BiCat(A, iY)$. Indeed, the optimal condition is to have $\Phi_{A,Y}$ be a strictly natural family of isomorphisms. Thus, we begin by noting that there are canonical icons 
 % https://q.uiver.app/#q=WzAsNSxbMSwyLCJcXEJpQ2F0KEMsaVkpIl0sWzAsMSwiXFxCaUNhdChCLGlZKSJdLFsxLDAsIlxcQmlDYXQoQSxpWSkiXSxbMywyLCJcXEJpQ2F0KEEsaVkpIl0sWzMsMCwiXFxCaUNhdChBLGlZKSJdLFswLDEsIkdeKiJdLFsxLDIsIkZeKiJdLFswLDIsIihGRyleKiIsMl0sWzMsNCwiMV97QX1eKiIsMix7ImN1cnZlIjozfV0sWzMsNCwiMSIsMCx7ImN1cnZlIjotM31dLFsxLDcsIiIsMix7InNob3J0ZW4iOnsidGFyZ2V0IjoyMH19XSxbOSw4LCIiLDAseyJzaG9ydGVuIjp7InNvdXJjZSI6MjAsInRhcmdldCI6MjB9fV1d
 \[\begin{tikzcd}[ampersand replacement=\&]
 	\& {\BiCat(A,iY)} \&\& {\BiCat(A,iY)} \\
 	{\BiCat(B,iY)} \\
 	\& {\BiCat(C,iY)} \&\& {\BiCat(A,iY)}
 	\arrow["{F^*}", from=2-1, to=1-2]
 	\arrow[""{name=0, anchor=center, inner sep=0}, "{(FG)^*}"', from=3-2, to=1-2]
 	\arrow["{G^*}", from=3-2, to=2-1]
 	\arrow[""{name=1, anchor=center, inner sep=0}, "{1_{A}^*}"', curve={height=18pt}, from=3-4, to=1-4]
 	\arrow[""{name=2, anchor=center, inner sep=0}, "1", curve={height=-18pt}, from=3-4, to=1-4]
 	\arrow[between={0}{0.8}, Rightarrow, nfold, from=2-1, to=0]
 	\arrow[between={0.2}{0.8}, Rightarrow, nfold, from=2, to=1]
 \end{tikzcd}\cdpunct[8pt]{.}\]
 If we assume our family of $2$-isomorphisms $\Phi_{A,Y}$ is strictly natural, this yields similar icons on $\Gray(\st(-),Y)$ by transporting along $\Phi$. Finally, within the context of $\Gray$-enriched Yoneda lemma, this naturally motivates specific definitions of the identity and compositor adjoint equivalences (or icons). The identity adjoint equivalence is relatively straight forward. 
 
 \begin{definition}
 	The identity adjoint equivalence $\mathbf{\iota}$ is defined by the following composite of pseudo natural transformation:
 	% https://q.uiver.app/#q=WzAsMTAsWzAsMCwiKiJdLFswLDQsIlxcQmlDYXQoQSxBKSJdLFsyLDIsIlxcQmlDYXQoQSxpIFxcc3QgQSkiXSxbNCwyLCJcXEdyYXkoXFxzdCBBLCBcXHN0IEEpIl0sWzAsNiwiXFxpb3RhIDoiXSxbMSw2LCIxX3tcXHN0IEF9Il0sWzEsNSwiXFxQaGleey0xfXNfe0F9Il0sWzMsNSwiXFxQaGleey0xfSgoU19BKV97Kn0xX0EpIl0sWzQsNSwiKFxcUGhpXnstMX0oU19BKV97Kn0pMV9BIl0sWzQsNiwiXFxzdF97QSxBfSAxX0EiXSxbMCwyLCJzX3tBfSIsMV0sWzAsMywiMV97XFxzdCBBfSJdLFsyLDMsIlxcUGhpXnstMX0iLDFdLFswLDEsIjFfQSIsMl0sWzEsMiwiKHNfe0F9KV97Kn0iLDFdLFsxLDMsIlxcc3Rfe0EsQX0iLDJdLFs2LDcsIjEgXFxzdGFyIFxcbWF0aHJte0NvaH1fe1xcQmlDYXR9IiwwLHsibGV2ZWwiOjJ9XSxbNSw2LCIiLDIseyJsZXZlbCI6Miwic3R5bGUiOnsiaGVhZCI6eyJuYW1lIjoibm9uZSJ9fX1dLFs3LDgsImFeey59IiwwLHsibGV2ZWwiOjIsInN0eWxlIjp7ImhlYWQiOnsibmFtZSI6Im5vbmUifX19XSxbOCw5LCIiLDAseyJsZXZlbCI6Miwic3R5bGUiOnsiaGVhZCI6eyJuYW1lIjoibm9uZSJ9fX1dLFs1LDksIiIsMSx7ImxldmVsIjoyfV0sWzEwLDEsIlxcbWF0aHJte0NvaH1fe1xcQmlDYXR9IiwxLHsic2hvcnRlbiI6eyJzb3VyY2UiOjIwLCJ0YXJnZXQiOjIwfX1dLFsxMSwyLCIxIiwyLHsic2hvcnRlbiI6eyJzb3VyY2UiOjIwLCJ0YXJnZXQiOjIwfX1dLFsyLDE1LCIxIiwyLHsic2hvcnRlbiI6eyJzb3VyY2UiOjIwLCJ0YXJnZXQiOjIwfX1dXQ==
 	\[\begin{tikzcd}[ampersand replacement=\&,column sep=small]
 		{*} \&\&\&\& \\
 		\\
 		\&\& {\BiCat(A,i \st A)} \&\& {\Gray(\st A, \st A)} \\
 		\\
 		{\BiCat(A,A)} \\
 		\& {\Phi^{-1}s_{A}} \&\& {\Phi^{-1}((S_A)_{*}1_A)} \& {(\Phi^{-1}(S_A)_{*})1_A} \\
 		{\iota :} \& {1_{\st A}} \&\&\& {\st_{A,A} 1_A}
 		\arrow[""{name=0, anchor=center, inner sep=0}, "{s_{A}}"{description}, from=1-1, to=3-3]
 		\arrow[""{name=1, anchor=center, inner sep=0}, "{1_{\st A}}", from=1-1, to=3-5]
 		\arrow["{1_A}"', from=1-1, to=5-1]
 		\arrow["{\Phi^{-1}}"{description}, from=3-3, to=3-5]
 		\arrow["{(s_{A})_{*}}"{description}, from=5-1, to=3-3]
 		\arrow[""{name=2, anchor=center, inner sep=0}, "{\st_{A,A}}"', from=5-1, to=3-5]
 		\arrow["{1 \star \mathrm{Coh}_{\BiCat}}", Rightarrow, nfold, from=6-2, to=6-4]
 		\arrow["{a^{.}}", equals, nfold, from=6-4, to=6-5]
 		\arrow[equals, nfold, from=6-5, to=7-5]
 		\arrow[equals, nfold, from=7-2, to=6-2]
 		\arrow[Rightarrow, nfold, from=7-2, to=7-5]
 		\arrow["1"', between={0.2}{0.8}, Rightarrow, nfold, from=1, to=3-3]
 		\arrow["{\mathrm{Coh}_{\BiCat}}"{description}, between={0.2}{0.8}, Rightarrow, nfold, from=0, to=5-1]
 		\arrow["1"', between={0.2}{0.8}, Rightarrow, nfold, from=3-3, to=2]
 	\end{tikzcd}\]
 \end{definition}

 We note that $\mathrm{Coh}_{\BiCat}$ is an invertible icon whose component at the unique $1$-cell $1_{*} \colon * \rightarrow *$ is depicted by the following diagram in $\BiCat(A, i \st A)$:
 % https://q.uiver.app/#q=WzAsNCxbMCwwLCJzX3tBfSJdLFszLDAsInNfe0F9Il0sWzAsMiwic197QX0iXSxbMywyLCJzX3tBfSJdLFswLDIsIiIsMCx7ImxldmVsIjoyLCJzdHlsZSI6eyJoZWFkIjp7Im5hbWUiOiJub25lIn19fV0sWzEsMywiIiwwLHsibGV2ZWwiOjIsInN0eWxlIjp7ImhlYWQiOnsibmFtZSI6Im5vbmUifX19XSxbMCwxLCIxX3tzX3tBfX0iLDAseyJsZXZlbCI6Mn1dLFsyLDMsIjFfe3Nfe0F9fSBcXHN0YXIgMV97QX0iLDIseyJsZXZlbCI6Mn1dLFs3LDYsIiIsMix7InNob3J0ZW4iOnsic291cmNlIjoyMCwidGFyZ2V0IjoyMH0sImxldmVsIjozfV1d
 \[\begin{tikzcd}[ampersand replacement=\&]
 	{s_{A}} \&\&\& {s_{A}} \\
 	\\
 	{s_{A}} \&\&\& {s_{A}}
 	\arrow[""{name=0, anchor=center, inner sep=0}, "{1_{s_{A}}}", Rightarrow, nfold, from=1-1, to=1-4]
 	\arrow[equals, nfold, from=1-1, to=3-1]
 	\arrow[equals, nfold, from=1-4, to=3-4]
 	\arrow[""{name=1, anchor=center, inner sep=0}, "{1_{s_{A}} \star 1_{A}}"', Rightarrow, nfold, from=3-1, to=3-4]
 	\arrow[between={0.2}{0.8}, Rightarrow, scaling nfold=3, from=1, to=0]
 \end{tikzcd}\cdpunct[8pt]{.}\]
 
 Thus, the composite invertible icon $1 \star \mathrm{Coh}_{\BiCat}$ is depicted in $\Gray(\st A, \st A)$ by
 % https://q.uiver.app/#q=WzAsNixbMCwwLCJcXHN0KDFfQSkiXSxbMywwLCJcXHN0KDFfQSkiXSxbMCwyLCJcXHN0KDFfQSkiXSxbMywyLCJcXHN0KDFfQSkiXSxbNSwwLCJhIl0sWzUsMiwiYSJdLFswLDIsIjFfe1xcc3QgMV9BfSIsMix7ImxldmVsIjoyLCJzdHlsZSI6eyJoZWFkIjp7Im5hbWUiOiJub25lIn19fV0sWzEsMywiMV97XFxzdCAxX0F9IiwwLHsibGV2ZWwiOjIsInN0eWxlIjp7ImhlYWQiOnsibmFtZSI6Im5vbmUifX19XSxbMCwxLCIxX3tcXHN0IDFfQX0iLDAseyJsZXZlbCI6Mn1dLFsyLDMsIlxcc3QoMV97MV9BfSkiLDIseyJsZXZlbCI6Mn1dLFs0LDUsIlxcdmFyZXBzaWxvbl97YX0iLDAseyJjdXJ2ZSI6LTN9XSxbNCw1LCIxX3thfSIsMix7ImN1cnZlIjozfV0sWzksOCwiXFxtYXRoYmZ7XFxpb3RhfSIsMix7InNob3J0ZW4iOnsic291cmNlIjoyMCwidGFyZ2V0IjoyMH0sImxldmVsIjozfV0sWzExLDEwLCJcXG1hdGhiZntcXGlvdGF9IiwyLHsic2hvcnRlbiI6eyJzb3VyY2UiOjIwLCJ0YXJnZXQiOjIwfX1dXQ==
 \[\begin{tikzcd}[ampersand replacement=\&,cramped]
 	{\st(1_A)} \&\&\& {\st(1_A)} \&\& a \\
 	\\
 	{\st(1_A)} \&\&\& {\st(1_A)} \&\& a
 	\arrow[""{name=0, anchor=center, inner sep=0}, "{1_{\st 1_A}}", Rightarrow, nfold, from=1-1, to=1-4]
 	\arrow["{1_{\st 1_A}}"', equals, nfold, from=1-1, to=3-1]
 	\arrow["{1_{\st 1_A}}", equals, nfold, from=1-4, to=3-4]
 	\arrow[""{name=1, anchor=center, inner sep=0}, "{\varepsilon_{a}}", curve={height=-18pt}, from=1-6, to=3-6]
 	\arrow[""{name=2, anchor=center, inner sep=0}, "{1_{a}}"', curve={height=18pt}, from=1-6, to=3-6]
 	\arrow[""{name=3, anchor=center, inner sep=0}, "{\st(1_{1_A})}"', Rightarrow, nfold, from=3-1, to=3-4]
 	\arrow["{\mathbf{\iota}}"', between={0.2}{0.8}, Rightarrow, nfold, from=2, to=1]
 	\arrow["{\mathbf{\iota}}"', between={0.2}{0.8}, Rightarrow, scaling nfold=3, from=3, to=0]
 \end{tikzcd}\cdpunct[8pt]{.}\]
 
 Unfortunately, the compositor adjoint equivalence $\mathbf{\chi}$ is for more complicated. 
 
 \begin{definition}
 	The compositor adjoint equivalence $\mathbf{\chi}$ is defined by the following composite of pseudonatural transformations:
 	% https://q.uiver.app/#q=WzAsOSxbMCwwLCJcXG1hdGhiZntcXGNoaX0gOiJdLFsxLDAsIlxcb3RpbWVzX3tcXEdyYXl9IFxcY2lyYyAoXFxzdCBcXHRpbWVzIFxcc3QpICJdLFszLDAsIlxcUGhpXnstMX0gXFxjaXJjIChzX3tBfSleeyp9IFxcY2lyYyBcXG90aW1lc197XFxCaUNhdH0gXFxjaXJjIChpIFxcdGltZXMgaSkgXFxjaXJjIChcXHN0IFxcdGltZXMgXFxzdCkiXSxbMywxLCJcXFBoaV57LTF9IFxcY2lyYyBcXG90aW1lc197XFxCaUNhdH0gXFxjaXJjICgxIFxcdGltZXMgKHNfe0F9KV57Kn0pIFxcY2lyYyAoaSBcXHRpbWVzIGkpIFxcY2lyYyAoXFxzdCBcXHRpbWVzIFxcc3QpIl0sWzMsMiwiXFxQaGleey0xfSBcXGNpcmMgXFxvdGltZXNfe1xcQmlDYXR9IFxcY2lyYyAoMSBcXHRpbWVzIChzX3tCfSlfeyp9KSBcXGNpcmMgKGkgXFx0aW1lcyAxKSBcXGNpcmMgKFxcc3QgXFx0aW1lcyAxKSJdLFszLDQsIlxcUGhpXnstMX0gXFxjaXJjIFxcb3RpbWVzX3tcXEJpQ2F0fSBcXGNpcmMgKChzX3tDfSlfeyp9IFxcdGltZXMgMSkgIl0sWzMsNSwiXFxQaGleey0xfSAgXFxjaXJjIChzX3tDfSlfeyp9IFxcY2lyYyBcXG90aW1lc197XFxCaUNhdH0iXSxbMywzLCJcXFBoaV57LTF9IFxcY2lyYyBcXG90aW1lc197XFxCaUNhdH0gIFxcY2lyYyAoKHNfe0J9KV57Kn0gXFx0aW1lcyAxKSAgXFxjaXJjIChpIFxcdGltZXMgMSkgXFxjaXJjIChcXHN0IFxcdGltZXMgMSkiXSxbMSw1LCJcXHN0IFxcY2lyYyBcXG90aW1lc197XFxCaUNhdH0iXSxbMSwyLCIiLDAseyJsZXZlbCI6Mn1dLFsyLDMsIiIsMCx7ImxldmVsIjoyfV0sWzMsNCwiIiwwLHsibGV2ZWwiOjJ9XSxbNSw2LCIiLDAseyJsZXZlbCI6Mn1dLFs0LDcsIiIsMCx7ImxldmVsIjoyfV0sWzcsNSwiIiwwLHsibGV2ZWwiOjJ9XSxbNiw4LCIiLDAseyJsZXZlbCI6Mn1dLFsxLDgsIiIsMix7ImxldmVsIjoyfV1d
 	\[\begin{tikzcd}[ampersand replacement=\&,cramped]
 		{\mathbf{\chi} \colon} \& {\otimes_{\Gray} \circ (\st \times \st) } \&\& {\Phi^{-1} \circ (s_{A})^{*} \circ \otimes_{\BiCat} \circ (i \times i) \circ (\st \times \st)} \\
 		\&\&\& {\Phi^{-1} \circ \otimes_{\BiCat} \circ (1 \times (s_{A})^{*}) \circ (i \times i) \circ (\st \times \st)} \\
 		\&\&\& {\Phi^{-1} \circ \otimes_{\BiCat} \circ (1 \times (s_{B})_{*}) \circ (i \times 1) \circ (\st \times 1)} \\
 		\&\&\& {\Phi^{-1} \circ \otimes_{\BiCat}  \circ ((s_{B})^{*} \times 1)  \circ (i \times 1) \circ (\st \times 1)} \\
 		\&\&\& {\Phi^{-1} \circ \otimes_{\BiCat} \circ ((s_{C})_{*} \times 1) } \\
 		\& {\st \circ \otimes_{\BiCat}} \&\& {\Phi^{-1}  \circ (s_{C})_{*} \circ \otimes_{\BiCat}}
 		\arrow[Rightarrow, nfold, from=1-2, to=1-4]
 		\arrow[Rightarrow, nfold, from=1-2, to=6-2]
 		\arrow[Rightarrow, nfold, from=1-4, to=2-4]
 		\arrow[Rightarrow, nfold, from=2-4, to=3-4]
 		\arrow[Rightarrow, nfold, from=3-4, to=4-4]
 		\arrow[Rightarrow, nfold, from=4-4, to=5-4]
 		\arrow[Rightarrow, nfold, from=5-4, to=6-4]
 		\arrow[Rightarrow, nfold, from=6-4, to=6-2]
 	\end{tikzcd}\]
 	Here, we have suppresed the associators or the unitors of $\BiCat$ for clarity. 
 \end{definition}
 
 %https://q.uiver.app/#q=WzAsMTEsWzAsMiwiXFxCaUNhdChCLEMpIFxcdGltZXMgXFxCaUNhdChBLEIpIl0sWzAsOCwiXFxCaUNhdChBLEMpIl0sWzMsNiwiXFxCaUNhdChBLGkgXFxjaXJjIFxcc3QgQykgIl0sWzUsOCwiXFxHcmF5KFxcc3QgQSwgXFxzdCBDKSJdLFsyLDAsIlxcR3JheShcXHN0IEIsXFxzdCBDKSBcXHRpbWVzIFxcR3JheShcXHN0IEEsIFxcc3QgQikiXSxbMywxLCJcXEdyYXkoXFxzdCBBLCBcXHN0IEMpIl0sWzMsNCwiXFxCaUNhdChpIFxcY2lyYyBcXHN0IEEsaSBcXGNpcmMgXFxzdCBDKSJdLFsyLDIsIlxcQmlDYXQoaSBcXGNpcmMgXFxzdCBCLGkgXFxjaXJjIFxcc3QgQykgXFx0aW1lcyBcXEJpQ2F0KGkgXFxjaXJjIFxcc3QgQSxpIFxcY2lyYyBcXHN0IEIpIl0sWzIsNCwiXFxCaUNhdChpIFxcY2lyYyBcXHN0IEIsaSBcXGNpcmMgXFxzdCBDKSBcXHRpbWVzIFxcQmlDYXQoaSBcXGNpcmMgXFxzdCBCKSJdLFsxLDYsIlxcQmlDYXQoQixpIFxcY2lyYyBcXHN0IEMpIFxcdGltZXMgXFxCaUNhdChBLEIpIl0sWzEsMywiXFxCaUNhdChpIFxcY2lyYyBcXHN0IEIsIGkgXFxjaXJjIFxcc3QgQykgXFx0aW1lcyBcXEJpQ2F0KEEsQikiXSxbMCwxXSxbMSwyLCIoc197Q30pX3sqfSIsMV0sWzIsMywiXFxQaGleey0xfSIsMV0sWzAsNCwiXFxzdCBcXHRpbWVzIFxcc3QiXSxbMSwzLCJcXHN0IiwyXSxbNCw1XSxbNSw2LCJpIiwxXSxbNiwyLCIoc197QX0pXnsqfSIsMV0sWzUsMywiMSJdLFs0LDcsImkgXFx0aW1lcyBpIiwxXSxbNyw2XSxbNyw4LCIxIFxcdGltZXMgKHNfe0F9KV57Kn0iLDFdLFs4LDJdLFswLDksIihzX3tDfSlfeyp9IFxcdGltZXMgMSIsMV0sWzksMl0sWzAsMTAsIihpIFxcY2lyYyBcXHN0KSBcXHRpbWVzIDEiLDFdLFsxMCw5LCIoc197Qn0pXnsqfSBcXHRpbWVzIDEiXSxbMTAsOCwiMSBcXHRpbWVzIChzX3tCfSlfeyp9Il1d

 
The action of $\mathbf{\chi}$ becomes more apparent when we note that it is also an invertible icon whose components are represented by the following chain of modifications in $\Gray(\st A, \st C)$:
% https://q.uiver.app/#q=WzAsMTMsWzEsNCwiXFxzdChcXGFscGhhX3tHfSBcXHN0YXIgXFxhbHBoYV97Rn0pIl0sWzAsMCwiXFxjaGlfe1xcYWxwaGFfe0d9LFxcYWxwaGFfe0Z9fSA6Il0sWzAsNCwiXFxQaGleey0xfSgxX3tzX3tjfX0gXFxzdGFyIChcXGFscGhhX3tHfSBcXHN0YXIgXFxhbHBoYV97Rn0pKSJdLFswLDMsIlxcUGhpXnstMX0oKDFfe3Nfe2N9fSBcXHN0YXIgXFxhbHBoYV97R30pIFxcc3RhciBcXGFscGhhX3tGfSkiXSxbMSwzLCJcXFBoaV57LTF9KChpXFxjaXJjIFxcc3QgXFxhbHBoYV97R30gXFxzdGFyIDFfe3Nfe0J9fSkgXFxzdGFyIFxcYWxwaGFfe0Z9KSJdLFsxLDIsIlxcUGhpXnstMX0oaVxcY2lyYyBcXHN0IFxcYWxwaGFfe0d9IFxcc3RhciAoMV97c197Qn19IFxcc3RhciBcXGFscGhhX3tGfSkpIl0sWzAsMiwiXFxQaGleey0xfShpXFxjaXJjIFxcc3QgXFxhbHBoYV97R30gXFxzdGFyIChpIFxcY2lyYyBcXHN0IFxcYWxwaGFfe0Z9IFxcc3RhciAxX3tzX3tBfX0pKSJdLFswLDEsIlxcUGhpXnstMX0oKGlcXGNpcmMgXFxzdCBcXGFscGhhX3tHfSBcXHN0YXIgaSBcXGNpcmMgXFxzdCBcXGFscGhhX3tGfSkgXFxzdGFyIDFfe3Nfe0F9fSkiXSxbMSwxLCJcXFBoaV57LTF9KChpXFxjaXJjIChcXHN0IFxcYWxwaGFfe0d9IFxcc3RhciBcXHN0IFxcYWxwaGFfe0Z9KSkgXFxzdGFyIDFfe3Nfe0F9fSkiXSxbMSwwLCJcXHN0KFxcYWxwaGFfe0d9KSBcXHN0YXIgXFxzdChcXGFscGhhX3tGfSkiXSxbMywwLCJHXzFGXzFhIl0sWzIsNCwiR18yRl8yYSJdLFszLDMsIkdfMkZfMWEiXSxbMCwyLCIiLDAseyJsZXZlbCI6Mywic3R5bGUiOnsiaGVhZCI6eyJuYW1lIjoibm9uZSJ9fX1dLFsyLDMsIiIsMCx7ImxldmVsIjozfV0sWzMsNCwiIiwwLHsibGV2ZWwiOjMsInN0eWxlIjp7ImhlYWQiOnsibmFtZSI6Im5vbmUifX19XSxbNCw1LCIiLDAseyJsZXZlbCI6M31dLFs1LDYsIiIsMCx7ImxldmVsIjozLCJzdHlsZSI6eyJoZWFkIjp7Im5hbWUiOiJub25lIn19fV0sWzYsNywiIiwwLHsibGV2ZWwiOjN9XSxbNyw4LCIiLDAseyJsZXZlbCI6Mywic3R5bGUiOnsiaGVhZCI6eyJuYW1lIjoibm9uZSJ9fX1dLFs4LDksIiIsMCx7ImxldmVsIjozLCJzdHlsZSI6eyJoZWFkIjp7Im5hbWUiOiJub25lIn19fV0sWzEwLDExLCJHXzIoXFxhbHBoYV97Rn0pX2EoXFxhbHBoYV97R30pX3tGXzFhfSAiLDJdLFsxMCwxMiwiKFxcYWxwaGFfe0d9KV97Rl8xYX0iXSxbMTIsMTEsIkdfMihcXGFscGhhX3tGfSlfYSJdLFsyMSwxMiwiXFxjaGkiLDAseyJzaG9ydGVuIjp7InNvdXJjZSI6MjAsInRhcmdldCI6MjB9fV1d
\[\begin{tikzcd}[ampersand replacement=\&,cramped,column sep=tiny]
	{\chi_{\alpha_{G},\alpha_{F}} :} \& {\st(\alpha_{G}) \star \st(\alpha_{F})} \&\& {G_1F_1a} \\
	{\Phi^{-1}((i\circ \st \alpha_{G} \star i \circ \st \alpha_{F}) \star 1_{s_{A}})} \& {\Phi^{-1}((i\circ (\st \alpha_{G} \star \st \alpha_{F})) \star 1_{s_{A}})} \\
	{\Phi^{-1}(i\circ \st \alpha_{G} \star (i \circ \st \alpha_{F} \star 1_{s_{A}}))} \& {\Phi^{-1}(i\circ \st \alpha_{G} \star (1_{s_{B}} \star \alpha_{F}))} \\
	{\Phi^{-1}((1_{s_{c}} \star \alpha_{G}) \star \alpha_{F})} \& {\Phi^{-1}((i\circ \st \alpha_{G} \star 1_{s_{B}}) \star \alpha_{F})} \&\& {G_2F_1a} \\
	{\Phi^{-1}(1_{s_{c}} \star (\alpha_{G} \star \alpha_{F}))} \& {\st(\alpha_{G} \star \alpha_{F})} \& {G_2F_2a}
	\arrow["{(\alpha_{G})_{F_1a}}", from=1-4, to=4-4]
	\arrow[""{name=0, anchor=center, inner sep=0}, "{G_2(\alpha_{F})_a(\alpha_{G})_{F_1a} }"', from=1-4, to=5-3]
	\arrow[Rightarrow, scaling nfold=3, no head, from=2-1, to=2-2]
	\arrow[Rightarrow, scaling nfold=3, no head, from=2-2, to=1-2]
	\arrow[Rightarrow, scaling nfold=3, from=3-1, to=2-1]
	\arrow[Rightarrow, scaling nfold=3, no head, from=3-2, to=3-1]
	\arrow[Rightarrow, scaling nfold=3, no head, from=4-1, to=4-2]
	\arrow[Rightarrow, scaling nfold=3, from=4-2, to=3-2]
	\arrow["{G_2(\alpha_{F})_a}", from=4-4, to=5-3]
	\arrow[Rightarrow, scaling nfold=3, from=5-1, to=4-1]
	\arrow[Rightarrow, scaling nfold=3, no head, from=5-2, to=5-1]
	\arrow["\chi", between={0.2}{0.8}, Rightarrow, nfold, from=0, to=4-4]
\end{tikzcd}\]

Both the tricategories $\BiCat$ and $\Gray$ can be considered $\BiCat_{2}$-enriched bicategories as deduced from \cite[Rem.~5.4]{gurski2013tricatcoh}, where $\BiCat_{2}$ is the strict $2$-category of bicategories, pseudofunctors, and Icons. Thus, the iconic structure of $\st$ implies that the trifunctor $\st$ can itself be considered as a $\BiCat_{2}$-enriched functor of $\BiCat_{2}$-enriched bicategories. In particular, any $3$-dimensional coherence data of the trifunctor $\st$ is merely induced by the coherence constraint cells. Hence, we have fully specified $\st$ as a trifunctor. 

\subsection{Encoding bicategorical pasting diagrams}

The primary usage of the strictification functor $\st$ is that it allows us to encode bicategorical pasting diagram in arbitrary bicategory $K$ using the canonical evaluation pseudofunctor $e\colon \st K \rightarrow K$. We motivate this with an example. 

Suppose we wish to specify the data of an adjunction in a bicategory $K$. Algebraically, it can be expressed using the following data:
% https://q.uiver.app/#q=WzAsMTYsWzEsMCwiZiA6IGEgXFxyaWdodGFycm93IGIiXSxbNCwwLCJnIDogYiBcXHJpZ2h0YXJyb3cgYSJdLFsxLDEsIlxcZXRhICA6IDFfe2F9IFxcUmlnaHRhcnJvdyBnZiJdLFs0LDEsIlxcZXBzaWxvbiAgOiBmZyBcXFJpZ2h0YXJyb3cgMV97Yn0iXSxbMCwzLCJmIl0sWzEsMiwiZjFfe2F9Il0sWzIsMiwiZihnZikiXSxbMywyLCIoZmcpZiJdLFs0LDIsIjFfe2J9ZiJdLFs1LDMsImYiXSxbMCw1LCJnIl0sWzUsNSwiZyJdLFsyLDQsIihnZilnIl0sWzMsNCwiZyhmZykiXSxbMSw0LCIxX3thfWciXSxbNCw0LCJnMV97Yn0iXSxbNCw1LCJyXnstMX0iLDAseyJsZXZlbCI6Mn1dLFs1LDYsIjFcXGV0YSIsMCx7ImxldmVsIjoyfV0sWzgsOSwibCIsMCx7ImxldmVsIjoyfV0sWzYsNywiYV57LTF9IiwwLHsibGV2ZWwiOjJ9XSxbNyw4LCJcXGVwc2lsb24xIiwwLHsibGV2ZWwiOjJ9XSxbNCw5LCIxX3tmfSIsMix7ImxldmVsIjoyfV0sWzE0LDEyLCJcXGV0YTEiLDAseyJsZXZlbCI6Mn1dLFsxMCwxNCwibF57LTF9IiwwLHsibGV2ZWwiOjJ9XSxbMTAsMTEsIjFfe2d9IiwyLHsibGV2ZWwiOjJ9XSxbMTMsMTUsIjFcXGVwc2lsb24iLDAseyJsZXZlbCI6Mn1dLFsxNSwxMSwiciIsMCx7ImxldmVsIjoyfV0sWzEyLDEzLCJhIiwwLHsibGV2ZWwiOjJ9XV0=
\[\begin{tikzcd}[ampersand replacement=\&,cramped]
	\& {f \colon a \rightarrow b} \&\&\& {g \colon b \rightarrow a} \& \\
	\& {\eta  \colon 1_{a} \Rightarrow gf} \&\&\& {\epsilon  \colon fg \Rightarrow 1_{b}} \\
	\& {f1_{a}} \& {f(gf)} \& {(fg)f} \& {1_{b}f} \\
	f \&\&\&\&\& f \\
	\& {1_{a}g} \& {(gf)g} \& {g(fg)} \& {g1_{b}} \\
	g \&\&\&\&\& g
	\arrow["{1\eta}", Rightarrow, nfold, from=3-2, to=3-3]
	\arrow["{a^{-1}}", Rightarrow, nfold, from=3-3, to=3-4]
	\arrow["\epsilon1", Rightarrow, nfold, from=3-4, to=3-5]
	\arrow["l", Rightarrow, nfold, from=3-5, to=4-6]
	\arrow["{r^{-1}}", Rightarrow, nfold, from=4-1, to=3-2]
	\arrow["{1_{f}}"', Rightarrow, nfold, from=4-1, to=4-6]
	\arrow["\eta1", Rightarrow, nfold, from=5-2, to=5-3]
	\arrow["a", Rightarrow, nfold, from=5-3, to=5-4]
	\arrow["{1\epsilon}", Rightarrow, nfold, from=5-4, to=5-5]
	\arrow["r", Rightarrow, nfold, from=5-5, to=6-6]
	\arrow["{l^{-1}}", Rightarrow, nfold, from=6-1, to=5-2]
	\arrow["{1_{g}}"', Rightarrow, nfold, from=6-1, to=6-6]
\end{tikzcd}\cdpunct[8pt]{.}\]
In a strict $2$-category such as $\nCat{1}$, we could express one of the triangle identity using the following pasting diagram:
% https://q.uiver.app/#q=WzAsOSxbMCwyLCJBIl0sWzEsMCwiQiJdLFsyLDIsIkEiXSxbMywwLCJCIl0sWzUsMiwiQiJdLFs2LDAsIkEiXSxbNywyLCJCIl0sWzgsMCwiQSJdLFs0LDEsIlxcZXF1aXYiXSxbMCwxLCJGIl0sWzEsMiwiRyIsMV0sWzIsMywiRiIsMl0sWzEsMywiMV97Qn0iXSxbMCwyLCIxX3tBfSIsMl0sWzQsNSwiRiJdLFs2LDcsIkYiLDJdLFs0LDYsIjFfe0J9IiwyXSxbNSw3LCIxX3tBfSJdLFs0LDcsIkYiLDJdLFsxMywxLCJcXGV0YSIsMSx7InNob3J0ZW4iOnsic291cmNlIjoyMCwidGFyZ2V0IjoyMH19XSxbMiwxMiwiXFxlcHNpbG9uIiwxLHsic2hvcnRlbiI6eyJzb3VyY2UiOjIwLCJ0YXJnZXQiOjIwfX1dXQ==
\[\begin{tikzcd}[ampersand replacement=\&,cramped]
	\& B \&\& B \&\&\& A \&\& A \\
	\&\&\&\& \equiv \\
	A \&\& A \&\&\& B \&\& B
	\arrow[""{name=0, anchor=center, inner sep=0}, "{1_{B}}", from=1-2, to=1-4]
	\arrow["G"{description}, from=1-2, to=3-3]
	\arrow["{1_{A}}", from=1-7, to=1-9]
	\arrow["F", from=3-1, to=1-2]
	\arrow[""{name=1, anchor=center, inner sep=0}, "{1_{A}}"', from=3-1, to=3-3]
	\arrow["F"', from=3-3, to=1-4]
	\arrow["F", from=3-6, to=1-7]
	\arrow["F"', from=3-6, to=1-9]
	\arrow["{1_{B}}"', from=3-6, to=3-8]
	\arrow["F"', from=3-8, to=1-9]
	\arrow["\eta"{description}, between={0.2}{0.8}, Rightarrow, nfold, from=1, to=1-2]
	\arrow["\epsilon"{description}, between={0.2}{0.8}, Rightarrow, nfold, from=3-3, to=0]
\end{tikzcd}\cdpunct[8pt]{.}\]
This motivates us to consider encoding such identities inside of $\st K$ instead of $K$ itself, where such identities are expressible using pasting diagram. To this end, we note the following relations under the evaluation map $e$ expressing both the cells $\eta$ and $\epsilon$ inside $\st K$:
% % https://q.uiver.app/#q=WzAsMTIsWzEsMCwiXFxzdCBLIl0sWzUsMCwiSyJdLFswLDIsImEiXSxbMSwxLCJiIl0sWzIsMiwiYSJdLFs0LDIsImEiXSxbNiwyLCJhIl0sWzAsNCwiYiJdLFsxLDMsImEiXSxbMiw0LCJiIl0sWzQsNCwiYiJdLFs2LDQsImIiXSxbMCwxLCJlIl0sWzIsMywiZiJdLFszLDQsImciXSxbMiw0LCJcXHZhcmVwc2lsb25fe2F9IiwyXSxbNSw2LCIxX3thfSIsMix7ImN1cnZlIjoyfV0sWzUsNiwiZ2YiLDAseyJjdXJ2ZSI6LTJ9XSxbNyw4LCJnIl0sWzgsOSwiZiJdLFs3LDksIlxcdmFyZXBzaWxvbl97Yn0iLDJdLFsxMCwxMSwiMV97Yn0iLDIseyJjdXJ2ZSI6Mn1dLFsxMCwxMSwiZmciLDAseyJjdXJ2ZSI6LTJ9XSxbMTUsMywiXFxldGEiLDIseyJzaG9ydGVuIjp7InNvdXJjZSI6MjAsInRhcmdldCI6MjB9fV0sWzE2LDE3LCJcXGV0YSIsMix7InNob3J0ZW4iOnsic291cmNlIjoyMCwidGFyZ2V0IjoyMH19XSxbOCwyMCwiXFxlcHNpbG9uIiwwLHsic2hvcnRlbiI6eyJzb3VyY2UiOjIwLCJ0YXJnZXQiOjIwfX1dLFsyMiwyMSwiXFxlcHNpbG9uIiwwLHsic2hvcnRlbiI6eyJzb3VyY2UiOjIwLCJ0YXJnZXQiOjIwfX1dXQ==
\[\begin{tikzcd}[ampersand replacement=\&,cramped]
	\& {\st K} \&\&\&\& K \& \\
	\& b \\
	a \&\& a \&\& a \&\& a \\
	\& a \\
	b \&\& b \&\& b \&\& b
	\arrow["e", from=1-2, to=1-6]
	\arrow["g", from=2-2, to=3-3]
	\arrow["f", from=3-1, to=2-2]
	\arrow[""{name=0, anchor=center, inner sep=0}, "{\varepsilon_{a}}"', from=3-1, to=3-3]
	\arrow[""{name=1, anchor=center, inner sep=0}, "{1_{a}}"', curve={height=12pt}, from=3-5, to=3-7]
	\arrow[""{name=2, anchor=center, inner sep=0}, "gf", curve={height=-12pt}, from=3-5, to=3-7]
	\arrow["f", from=4-2, to=5-3]
	\arrow["g", from=5-1, to=4-2]
	\arrow[""{name=3, anchor=center, inner sep=0}, "{\varepsilon_{b}}"', from=5-1, to=5-3]
	\arrow[""{name=4, anchor=center, inner sep=0}, "{1_{b}}"', curve={height=12pt}, from=5-5, to=5-7]
	\arrow[""{name=5, anchor=center, inner sep=0}, "fg", curve={height=-12pt}, from=5-5, to=5-7]
	\arrow["\eta"', between={0.2}{0.8}, Rightarrow, nfold, from=0, to=2-2]
	\arrow["\eta"', between={0.2}{0.8}, Rightarrow, nfold, from=1, to=2]
	\arrow["\epsilon", between={0.2}{0.8}, Rightarrow, nfold, from=4-2, to=3]
	\arrow["\epsilon", between={0.2}{0.8}, Rightarrow, nfold, from=5, to=4]
\end{tikzcd}\cdpunct[8pt]{.}\]

In general, let $\gamma$ be some $2$-cell between $e(\vec{f},b)$ and $e(\vec{g},b')$, where $b$ and $b'$ signify any arbitrary choice of bracketing in the evaluation. This $2$-cell can be sent inside $\st K$ using the embedding $s_{K} \colon K \rightarrow \st K$. We do this by mapping $\gamma \in K$ to the $2$-cell $\vec{f} \xRightarrow{\mathrm{Coh}} e(\vec{f},b) \xRightarrow{\gamma} e(\vec{g},b') \xRightarrow{\mathrm{Coh}^{-1}} \vec{g}$ in $\st K$, whenever $\vec{f}$ and $\vec{g}$ contains no identity $1$-cells. If they do, we also remove them using coherence.

With this idea, we now embed one of the triangle identity of $(f,g,\eta,\epsilon)$ inside $\st k$: 
% https://q.uiver.app/#q=WzAsOSxbMSwzLCJiIl0sWzIsMSwiYSJdLFszLDMsImIiXSxbMCwxLCJhIl0sWzQsMiwiXFxlcXVpdiJdLFs1LDEsImEiXSxbOCwzLCJiIl0sWzMsMCwiZiJdLFswLDQsImYiXSxbMCwxLCJnIiwxXSxbMSwyLCJmIiwxXSxbMCwyLCJcXHZhcmVwc2lsb25fe2J9IiwxXSxbMywwLCJmIiwxXSxbMywxLCJcXHZhcmVwc2lsb25fe2F9IiwxXSxbNSw2LCJmIiwwLHsiY3VydmUiOi0zfV0sWzUsNiwiZiIsMix7ImN1cnZlIjozfV0sWzMsNywiIiwwLHsic3R5bGUiOnsiaGVhZCI6eyJuYW1lIjoibm9uZSJ9fX1dLFs3LDJdLFszLDgsIiIsMCx7InN0eWxlIjp7ImhlYWQiOnsibmFtZSI6Im5vbmUifX19XSxbOCwyXSxbNywxLCIiLDAseyJzaG9ydGVuIjp7InNvdXJjZSI6MjAsInRhcmdldCI6MjB9LCJsZXZlbCI6Mn1dLFswLDgsIiIsMCx7InNob3J0ZW4iOnsic291cmNlIjoyMCwidGFyZ2V0IjoyMH0sImxldmVsIjoyfV0sWzEsMTEsIlxcZXBzaWxvbiIsMCx7InNob3J0ZW4iOnsic291cmNlIjoyMCwidGFyZ2V0IjoyMH19XSxbMTMsMCwiXFxldGEiLDAseyJzaG9ydGVuIjp7InNvdXJjZSI6MjAsInRhcmdldCI6MjB9fV0sWzE0LDE1LCIxX3tmfSIsMCx7InNob3J0ZW4iOnsic291cmNlIjoyMCwidGFyZ2V0IjoyMH19XV0=
\[\begin{tikzcd}[ampersand replacement=\&]
	\&\&\& f \&\&\&\&\& \\
	a \&\& a \&\&\& a \\
	\&\&\&\& \equiv \\
	\& b \&\& b \&\&\&\&\& b \\
	f
	\arrow[between={0.2}{0.8}, Rightarrow, nfold, from=1-4, to=2-3]
	\arrow[from=1-4, to=4-4]
	\arrow[no head, from=2-1, to=1-4]
	\arrow[""{name=0, anchor=center, inner sep=0}, "{\varepsilon_{a}}"{description}, from=2-1, to=2-3]
	\arrow["f"{description}, from=2-1, to=4-2]
	\arrow[no head, from=2-1, to=5-1]
	\arrow["f"{description}, from=2-3, to=4-4]
	\arrow[""{name=1, anchor=center, inner sep=0}, "f", curve={height=-18pt}, from=2-6, to=4-9]
	\arrow[""{name=2, anchor=center, inner sep=0}, "f"', curve={height=18pt}, from=2-6, to=4-9]
	\arrow["g"{description}, from=4-2, to=2-3]
	\arrow[""{name=3, anchor=center, inner sep=0}, "{\varepsilon_{b}}"{description}, from=4-2, to=4-4]
	\arrow[between={0.2}{0.8}, Rightarrow, nfold, from=4-2, to=5-1]
	\arrow[from=5-1, to=4-4]
	\arrow["\eta", between={0.2}{0.8}, Rightarrow, nfold, from=0, to=4-2]
	\arrow["\epsilon", between={0.2}{0.8}, Rightarrow, nfold, from=2-3, to=3]
	\arrow["{1_{f}}", between={0.2}{0.8}, Rightarrow, nfold, from=1, to=2]
\end{tikzcd}\cdpunct[8pt]{.}\]
Of course, $\varepsilon_{a}$ and $\varepsilon_{b}$ are strict identities inside $\st K$. Thus, we could equivalently write
% https://q.uiver.app/#q=WzAsOSxbMSwyLCJiIl0sWzIsMCwiYSJdLFszLDIsImIiXSxbMCwwLCJhIl0sWzQsMSwiXFxlcXVpdiJdLFs1LDAsImEiXSxbOCwyLCJiIl0sWzYsMiwiYiJdLFs3LDAsImEiXSxbMCwxLCJnIiwxXSxbMSwyLCJmIl0sWzAsMiwiXFx2YXJlcHNpbG9uX3tifSIsMl0sWzMsMCwiZiIsMl0sWzMsMSwiXFx2YXJlcHNpbG9uX3thfSJdLFs1LDYsImYiLDFdLFs1LDcsImYiLDJdLFs3LDYsIlxcdmFyZXBzaWxvbl97Yn0iLDJdLFs4LDYsImYiXSxbNSw4LCJcXHZhcmVwc2lsb25fe2F9Il0sWzEsMTEsIlxcZXBzaWxvbiIsMCx7InNob3J0ZW4iOnsic291cmNlIjoyMCwidGFyZ2V0IjoyMH19XSxbMTMsMCwiXFxldGEiLDAseyJzaG9ydGVuIjp7InNvdXJjZSI6MjAsInRhcmdldCI6MjB9fV0sWzgsMTQsIiIsMCx7InNob3J0ZW4iOnsic291cmNlIjoyMCwidGFyZ2V0IjoyMH19XSxbMTQsNywiIiwwLHsic2hvcnRlbiI6eyJzb3VyY2UiOjIwLCJ0YXJnZXQiOjIwfX1dXQ==
\[\begin{tikzcd}[ampersand replacement=\&]
	a \&\& a \&\&\& a \&\& a \& \\
	\&\&\&\& \equiv \\
	\& b \&\& b \&\&\& b \&\& b
	\arrow[""{name=0, anchor=center, inner sep=0}, "{\varepsilon_{a}}", from=1-1, to=1-3]
	\arrow["f"', from=1-1, to=3-2]
	\arrow["f", from=1-3, to=3-4]
	\arrow["{\varepsilon_{a}}", from=1-6, to=1-8]
	\arrow["f"', from=1-6, to=3-7]
	\arrow[""{name=1, anchor=center, inner sep=0}, "f"{description}, from=1-6, to=3-9]
	\arrow["f", from=1-8, to=3-9]
	\arrow["g"{description}, from=3-2, to=1-3]
	\arrow[""{name=2, anchor=center, inner sep=0}, "{\varepsilon_{b}}"', from=3-2, to=3-4]
	\arrow["{\varepsilon_{b}}"', from=3-7, to=3-9]
	\arrow["\eta", between={0.2}{0.8}, Rightarrow, nfold, from=0, to=3-2]
	\arrow["\epsilon", between={0.2}{0.8}, Rightarrow, nfold, from=1-3, to=2]
	\arrow[between={0.2}{0.8}, Rightarrow, nfold, from=1, to=3-7]
	\arrow[between={0.2}{0.8}, Rightarrow, nfold, from=1-8, to=1]
\end{tikzcd}\cdpunct[8pt]{}\]
to express the same triangle identity. Finally, Using the evaluation map $e \colon \st B \rightarrow B$, it is relatively straightforward to verify that this pasting diagram equality correspond to the following equality of $2$-cells inside $B$:
% https://q.uiver.app/#q=WzAsOSxbMCwxLCJmIl0sWzEsMCwiZjFfe2F9Il0sWzIsMCwiZihnZikiXSxbNCwwLCIoZmcpZiJdLFs1LDAsIjFfe2J9ZiJdLFs2LDEsImYiXSxbMCwyLCJmIl0sWzYsMiwiZiJdLFszLDEsIlxcdmVxdWl2Il0sWzAsMSwicl57LTF9IiwwLHsibGV2ZWwiOjJ9XSxbMSwyLCIxXFxldGEiLDAseyJsZXZlbCI6Mn1dLFs0LDUsImwiLDAseyJsZXZlbCI6Mn1dLFsyLDMsImFeey0xfSIsMCx7ImxldmVsIjoyfV0sWzMsNCwiXFxlcHNpbG9uMSIsMCx7ImxldmVsIjoyfV0sWzYsNywiMV97Zn0iLDIseyJsZXZlbCI6Mn1dXQ==
\[\begin{tikzcd}[ampersand replacement=\&]
	\& {f1_{a}} \& {f(gf)} \&\& {(fg)f} \& {1_{b}f} \& \\
	f \&\&\& \vequiv \&\&\& f \\
	f \&\&\&\&\&\& f
	\arrow["{1\eta}", Rightarrow, nfold, from=1-2, to=1-3]
	\arrow["{a^{-1}}", Rightarrow, nfold, from=1-3, to=1-5]
	\arrow["\epsilon1", Rightarrow, nfold, from=1-5, to=1-6]
	\arrow["l", Rightarrow, nfold, from=1-6, to=2-7]
	\arrow["{r^{-1}}", Rightarrow, nfold, from=2-1, to=1-2]
	\arrow["{1_{f}}"', Rightarrow, nfold, from=3-1, to=3-7]
\end{tikzcd}\cdpunct[8pt]{.}\]
Therefore, as a general note, whenever we write bicategorical pasting diagram inside a given bicategory $K$, we always mean to read such diagrams inside $\st K$ with the use of the evaluation functor $e \colon\st K \rightarrow K$. A more formal treatment of this idea is given using pasting schemes. See \cite[Ch.~3]{yau2cat2021} for more details. 